% Date : 04/04/2023
% Version : 3
% Auteur : SBA
% Relu : oui
% Relecteur : Chahira Hajlaoui, Colas Bardavid, Jimmy Roussel/ALD
% Harmonisateur: CBN, GBU

% ------------------------------------------------------------ %
% ----------------------  Fiche ELC04  ----------------------- %
% Thème : Énergie et puissance électriques
% Classes : Toutes
% ------------------------------------------------------------ %



% ======================================================= % 
% ============== Gestion de la compilation ============== %
% ======================================================= % 
\ifdefined\mainIsLoaded\else
\RequirePackage{import}
\subimport{../../}{_preambule_CdE_PC}
\begin{document}
\fi
% ======================================================= % 
% ======================================================= % 



% ======================================================= % 
% -----------------------------------%
% La commande \couleurNBouCouleur prend deux paramètres :
% #1 est la couleur NB ; #2 est la couleur "en couleur"
\def\JRLblue{\couleurNBouCouleur{black}{blue!75}}
 \def\JRLblueC{\couleurNBouCouleur{lightgray!30}{blue!15}}
 \def\JRLblueF{\couleurNBouCouleur{gray}{blue!50}}
\def\JRLred{\couleurNBouCouleur{gray}{red!66}}
 \def\JRLcyan{\couleurNBouCouleur{gray}{cyan}}
 \def\JRLcyanclair{\couleurNBouCouleur{lightgray!30}{cyan!10}}
% ======================================================= % 





% ============================================================================ %
%                            FICHE D'ENTRAÎNEMENT                              %
% ============================================================================ %
% ======================= Métadonnées sur la fiche =========================== %
\titreFicheEntrainement		{Énergie et puissance électriques}
\grandTheme					{Électricité}
\numeroFiche				{ELC04}
\uniqueID					{JzWYjOTLDiI} % une chaîne de caractères (a-z, A-Z) aléatoire de 10 caractères.
\nombreColonnesReponses		{2} % 1, 2, 3, etc.
% ============================================================================ %

%%%%%%%%%%%%%%%%%%%%%%%%%%%%%%%%%%%%%%%%%%%%%%%%%%%%%%%%%%%%%%%%%%%%%%%%%%%%%%%%%%%%%%%%%%%%%%
\debutFicheEntrainement                                                                      %
%%%%%%%%%%%%%%%%%%%%%%%%%%%%%%%%%%%%%%%%%%%%%%%%%%%%%%%%%%%%%%%%%%%%%%%%%%%%%%%%%%%%%%%%%%%%%%




% --------------------------------------------------------------------- %
%                              Prérequis                                %
%                             (facultatif)                              %
% --------------------------------------------------------------------- %
% ================== Métadonnées sur le prérequis ===================== %
\avecPrerequis{Y} % Y ou N
% ===================================================================== %
\begin{prerequis}
	Puissance électrique. Relation puissance-énergie. Expressions des énergies stockées dans une bobine et dans un condensateur. Effet Joule.
\end{prerequis}
% --------------------------------------------------------------------- %






% •••••••••••••••••••••••••••••••••••••••••••••••••••••••••••••••••••••••••••• %
% •••••••••••••••••••••••••••••••••••••••••••••••••••••••••••••••••••••••••••• %
\sectionFicheEntrainement{Pour commencer}
% •••••••••••••••••••••••••••••••••••••••••••••••••••••••••••••••••••••••••••• %
% •••••••••••••••••••••••••••••••••••••••••••••••••••••••••••••••••••••••••••• %






% ********************************************************************* %
%                           ENTRAÎNEMENT                                % 
% ********************************************************************* %
% ================ Métadonnées sur l'entraînement ===================== %

\titreEntrainementFacultatif	{Puissance et énergie}
\hauteurLargeurCadreReponse		{8mm}{2.5cm}
\dureeResolutionFacultative		{1} % 1, 2, 3 ou 4
\typeEntrainement				{AN} % Newton ou QCM ou AN ou vide (exercice "normal")
\basiqueEtTransversal			{Y}
\nombreColonnesQuestions		{1} % vide, 1, 2, 3, etc.
\avecPlusieursQuestions			{Y} % Y ou N
\initialisationEntrainement
% ===================================================================== %




% ^^^^^^^^^^^^^^^^^^^^^^^^^^^^^^^^^^^^^^ %
%               Question                 %
% ^^^^^^^^^^^^^^^^^^^^^^^^^^^^^^^^^^^^^^ %
Le chargeur %de la batterie 
d'un téléphone portable consomme une puissance %moyenne
 de \SI{5}{\watt}. La charge complète de la batterie (à partir d'une batterie vide)
  prend 
\SI{55}{\min}. 

\bigskip

Calculer l'énergie $E$ contenue dans la batterie:

\minusSmallskip

% =============================================================================== %
\debutEntrainement
% =============================================================================== %

% ^^^^^^^^^^^^^^^^^^^^^^^^^^^^^^^^^^^^^^ %
%               Question                 %
% ^^^^^^^^^^^^^^^^^^^^^^^^^^^^^^^^^^^^^^ %

\begin{enonce}
en joules
\end{enonce}

\reponse{\SI{16,5}{\kJ}}

\begin{corrige}
L'énergie contenue dans la batterie vaut $E=P\Delta t$ où $P=\SI{5}{\watt}$ et $\Delta t=\SI{55}{\min}=\num{55 }\times \SI{60}{s}=\SI{3300}{\second}$. L'énergie vaut donc
$E=5 \times \num{3300} \, \si{\joule}=\SI{16,5}{\kJ}$.
\end{corrige}

% ************************************** %


% ^^^^^^^^^^^^^^^^^^^^^^^^^^^^^^^^^^^^^^ %
%               Question                 %
% ^^^^^^^^^^^^^^^^^^^^^^^^^^^^^^^^^^^^^^ %
\begin{enonce}
en watt-heures (\si{\watt.\hour})
\end{enonce}

\reponse{\SI{4,6}{\watt.\hour}}

\begin{corrige}
L'énergie contenue dans la batterie vaut $E=\SI{16,5}{\kJ}$. Par ailleurs, $e=\SI{1}{\watt.\hour}$ est l'énergie consommée à une puissance de $\SI{1}{\watt}$ pendant $\SI{1}{\hour}$, soit $e=\SI{1}{W}\times \SI{3600}{s} =\SI{3,6}{\kJ}$.

On a donc $E = \frac{\SI{16,5}{\kJ}}{\SI{3,6}{\kJ}} \times \SI{1}{\watt.\hour} = \SI{4,6}{\watt.\hour}$.
\end{corrige}

% ************************************** %

% =============================================================================== %
\finEntrainement
% =============================================================================== %




% ********************************************************************* %
%                           ENTRAÎNEMENT                                % 
% ********************************************************************* %
% ================ Métadonnées sur l'entraînement ===================== %

\titreEntrainementFacultatif	{Voiture de série contre Formule 1}
\hauteurLargeurCadreReponse		{8mm}{4cm}
\dureeResolutionFacultative		{1} % 1, 2, 3 ou 4
\typeEntrainement				{AN} % Newton ou QCM ou AN ou vide (exercice "normal")
\basiqueEtTransversal			{Y}
\nombreColonnesQuestions		{1} % vide, 1, 2, 3, etc.
\avecPlusieursQuestions			{Y} % Y ou N
\initialisationEntrainement
% ===================================================================== %



% =============================================================================== %
\debutEntrainement
% =============================================================================== %


Les voitures de courses \glm{Formule 1} sont des véhicules hybrides : elles possèdent à la fois un moteur thermique 
et un moteur électrique. On souhaite comparer le moteur électrique d'une Formule 1 à celui d'une simple voiture de série.

On donne les informations suivantes:

\begin{center}
\def\monEspace{\mbox{\phantom{aa}}}
\renewcommand{\arraystretch}{1.2}
\begin{tabular}{|c|c|c|}
\hline
& \monEspace{}\textbf{Hyundai Ioniq 6}\monEspace{} & \monEspace{}\textbf{Formule 1}\monEspace{} \\
\hline
\textbf{Capacité batterie} & \SI{77,4}{\kilo\watt . \hour} & \SI{4}{\mega\joule}\\
\hline
\textbf{Puissance moteur} &\SI{239}{\kilo\watt} & \SI{160}{cv} \\
\hline
\monEspace{}\textbf{Consommation moyenne}\monEspace{} & \SI{15,1}{\kWh \per 100 \kilo\meter} & \\
\hline
\end{tabular}
\renewcommand{\arraystretch}{1}
\end{center}

On indique que $\SI{1}{cv}=\SI{0,735}{\kW}$.

% ^^^^^^^^^^^^^^^^^^^^^^^^^^^^^^^^^^^^^^ %
%               Question                 %
% ^^^^^^^^^^^^^^^^^^^^^^^^^^^^^^^^^^^^^^ %

\begin{enonce}
Calculer l'autonomie en \si{\kilo\meter} de la batterie de la Hyundai Ioniq 6
\end{enonce}

\reponse{\SI{513}{\kilo\meter}}

\begin{corrige}
L'énergie contenue dans la batterie vaut $E=\SI{77,4}{\kilo\watt.\hour}$. 

La consommation moyenne valant $C=\SI{15,1}{\kWh \per 100 \kilo\meter}$, l'autonomie en kilomètres vaut
\[
\frac{E}{C}
=
\frac{\SI{77,4}{\kilo\watt .\hour}}{\SI{15,1}{\kWh \per 100 \kilo\meter}}=\SI{513}{\km}.
\]
\end{corrige}

% ************************************** %


% ^^^^^^^^^^^^^^^^^^^^^^^^^^^^^^^^^^^^^^ %
%               Question                 %
% ^^^^^^^^^^^^^^^^^^^^^^^^^^^^^^^^^^^^^^ %

\begin{enonce}
Quel véhicule possède la batterie de plus grande capacité ?
\end{enonce}

\reponse{Hyundai Ioniq 6} 

\begin{corrige}
En reprenant le calcul de la question précédente, $e=\SI{1}{ \watt / \hour}=\SI{3,6}{\kJ}$, donc l'énergie totale stockée dans les batteries des voitures de série vaut en Joule 
$E=\num{77,4 e3}\times \num{3,6 e3}\, \si{J}=\SI{279}{\mega \joule}$. C'est donc la voiture de série qui possède la batterie de plus grande capacité.
\end{corrige}

% ************************************** %


% ^^^^^^^^^^^^^^^^^^^^^^^^^^^^^^^^^^^^^^ %
%               Question                 %
% ^^^^^^^^^^^^^^^^^^^^^^^^^^^^^^^^^^^^^^ %

\begin{enonce}
Quel véhicule possède le moteur électrique le plus puissant ?
\end{enonce}

\reponse{Hyundai Ioniq 6} 

\begin{corrige}
La puissance en \si{cv} du moteur de la voiture électrique de série vaut $\mathcal{P}=\num{239} /\num{0,735} \, \si{cv}=\SI{325}{cv}$.
\end{corrige}

% ************************************** %

% =============================================================================== %
\finEntrainement
% =============================================================================== %



\newpage



% ********************************************************************* %
%                           ENTRAÎNEMENT                                % 
% ********************************************************************* %
% ================ Métadonnées sur l'entraînement ===================== %

\titreEntrainementFacultatif	{Identifications de courbes}
\hauteurLargeurCadreReponse		{8mm}{4cm}
\dureeResolutionFacultative		{4} % 1, 2, 3 ou 4
\typeEntrainement				{} % Newton ou QCM ou AN ou vide (exercice "normal")
\basiqueEtTransversal			{Y}
\nombreColonnesQuestions		{1} % vide, 1, 2, 3, etc.
\avecPlusieursQuestions			{Y} % Y ou N
\initialisationEntrainement
% ===================================================================== %

\def\globalScaleGraphes{0.7}

% =============================================================================== %
\debutEntrainement
% =============================================================================== %

Une tension  $u(t)$  est appliquée aux bornes d'un conducteur ohmique de résistance $R=\SI{10}{\Omega}$. 

Identifier parmi les courbes proposées celle correspondant à la puissance 
\[
\mathcal{P}(t)=\dfrac{u^2(t)}{R}
\]
 dissipée par effet Joule dans la résistance.

\begin{enonce}
Pour $u(t)=3\sin(\omega t)$ avec $\omega = 2\pi\; \si{rad.s^{-1}}$.

\begin{listeQCM2Colonnes}
\item \subimport{_images/}{fiche_ELC04_puissance_Q1a_CBD_v2.tex}
\item \subimport{_images/}{fiche_ELC04_puissance_Q1b_CBD_v2.tex}
\item \subimport{_images/}{fiche_ELC04_puissance_Q1c_CBD_v2.tex}
\item \subimport{_images/}{fiche_ELC04_puissance_Q1d_CBD_v2.tex}
\end{listeQCM2Colonnes}
\medskip
\end{enonce}

\reponse{\reponseA{}} 

\begin{corrige}
La puissance reçue par la résistance s'écrit $\mathcal{P}=\dfrac{u^2}{R}$. Ici, on a donc
\[
\mathcal{P}=\dfrac{9}{10}\sin^2 \left( \omega t\right)=\dfrac{9}{20}\big ( 1-\cos \left( 2 \omega t\right)\big ).
\] 
La puissance a donc une valeur moyenne de $\dfrac{9}{20}$, une valeur maximale de $\dfrac{9}{10}$ et une période $T=\SI{0.5}{s}$.

C'est la réponse \reponseA{} qui est la bonne.
\end{corrige}

% ************************************** %


% ^^^^^^^^^^^^^^^^^^^^^^^^^^^^^^^^^^^^^^ %
%               Question                 %
% ^^^^^^^^^^^^^^^^^^^^^^^^^^^^^^^^^^^^^^ %

\begin{enonce}
Pour $u(t)=1+2\cos(\omega t)$ avec $\omega = \pi \;\si{rad.s^{-1}}$.

\begin{listeQCM2Colonnes}
\item \subimport{_images/}{fiche_ELC04_puissance_Q2c_CBD_v2.tex}
\item \subimport{_images/}{fiche_ELC04_puissance_Q2b_CBD_v2.tex}
\item \subimport{_images/}{fiche_ELC04_puissance_Q2a_CBD_v2.tex}
\item \subimport{_images/}{fiche_ELC04_puissance_Q2d_CBD_v2.tex}
\end{listeQCM2Colonnes}
\medskip
\end{enonce}

\reponse{\reponseC{}} 

\begin{corrige}
Commençons par linéariser l'expression de la puissance. On a 
\begin{align*}
\mathcal{P}(t) & =\dfrac{u^{2}}{R}=\dfrac{1}{10}\big(1+2\cos(\omega t)\big)^{2}=\dfrac{1}{10}\left(1+4\cos^{2}(\omega t)+4\cos(\omega t)\right)\\
 & =\dfrac{1}{10}\big(3+2\cos(2\omega t)+4\cos(\omega t)\big).
\end{align*}
On constate que la puissance est maximale à $t=0$. De plus, la composante fondamentale
de ce signal est de période égale à $T_{\text{fondamental}}=\dfrac{2\pi}{\omega}=\SI{2}{s}$.
Finalement, comme $u(t)$ s'annule (par exemple en $\omega t=\dfrac{\pi}{3}$),
la puissance s'annule aussi. 

Il n'y a qu'une courbe qui vérifie ces conditions: c'est la \reponseC{} qui est la bonne.
\end{corrige}

% ************************************** %


\newpage


% ^^^^^^^^^^^^^^^^^^^^^^^^^^^^^^^^^^^^^^ %
%               Question                 %
% ^^^^^^^^^^^^^^^^^^^^^^^^^^^^^^^^^^^^^^ %

\begin{enonce}
Pour $u(t)=3\exp \left(-\dfrac{t}{\tau}\right)$ avec $\tau=\SI{2}{s}$.\\

\begin{listeQCM2Colonnes}
\item \subimport{_images/}{fiche_ELC04_puissance_Q3b_CBD_v2.tex}
\item \subimport{_images/}{fiche_ELC04_puissance_Q3c_CBD_v2.tex}
\item \subimport{_images/}{fiche_ELC04_puissance_Q3a_CBD_v2.tex}
\end{listeQCM2Colonnes}
\medskip
\end{enonce}

\reponse{\reponseC{}} 

\begin{corrige}
La puissance a pour expression $\mathcal{P}=\dfrac{u^{2}}{R}=\dfrac{9}{10}\exp\left(-\dfrac{2t}{\tau}\right)$. On a donc
\[
\diff{\mathcal{P}(t)}{t} = -\frac{2}{\tau}\dfrac{9}{10}\exp\left(-\dfrac{2t}{\tau}\right)
\qquad\text{donc}\qquad
\diff{\mathcal{P}(t)}{t}(t=0) = -\frac{2}{\tau}\dfrac{9}{10} = - \dfrac{9}{10} \si{W.s^{-1}}.
\] 
En exploitant la pente à l'origine, on trouve que c'est la réponse \reponseC{} qui est la bonne.
\end{corrige}

% ************************************** %


% =============================================================================== %
\finEntrainement
% =============================================================================== %













% ********************************************************************* %
%                           ENTRAÎNEMENT                                % 
% ********************************************************************* %
% ================ Métadonnées sur l'entraînement ===================== %

\pointInterrogation				{N}
\titreEntrainementFacultatif	{Un calcul graphique}
\hauteurLargeurCadreReponse		{8mm}{3.5cm}
\dureeResolutionFacultative		{1} % 1, 2, 3 ou 4
\typeEntrainement				{} % Newton ou QCM ou AN ou vide (exercice "normal")
\basiqueEtTransversal			{Y}
\nombreColonnesQuestions		{1} % vide, 1, 2, 3, etc.
\avecPlusieursQuestions			{N} % Y ou N
\initialisationEntrainement
% ===================================================================== %



Pour un dipôle soumis à un signal alternatif harmonique, la puissance moyenne vaut:
\[
\mathcal{P}_{moy}=\dfrac{U_{0}I_{0}}{2}\cos(\varphi)
\]
où $U_{0}$ et $I_{0}$ sont respectivement l'amplitude de la tension et du courant et où $\varphi$ représente la valeur du déphasage angulaire entre la tension et l'intensité du courant. 

La figure ci-dessous donne les représentations graphiques de la tension $u(t)$ et de l'intensité $i(t)$ en convention récepteur. 

\begin{center}
\subimport{_images/}{fiche_ELC04_dephasage_CBD_v2}
\end{center}

% =============================================================================== %
\debutEntrainement
% =============================================================================== %


% ^^^^^^^^^^^^^^^^^^^^^^^^^^^^^^^^^^^^^^ %
%               Question                 %
% ^^^^^^^^^^^^^^^^^^^^^^^^^^^^^^^^^^^^^^ %

\begin{enonce}
Déterminer la puissance moyenne reçue par ce dipôle
\end{enonce}

\reponse{$\SI{3.75}{W}$}

\begin{corrige}
On lit graphiquement une période de $T=\SI{3}{ms}$ et un décalage temporel $\Delta t=\SI{0.5}{ms}$ entre les deux signaux. Le déphasage est donc $\varphi=2\pi\dfrac{\Delta t}{T}=\dfrac{\pi}{3}\;\si{rad}$. Donc, $\cos(\varphi) = \frac{1}{2}$.

Les amplitudes de la tension et de l'intensité sont respectivement $U_{0}=\SI{3}{V}$ et $I_{0}=\SI{5}{A}$. La puissance moyenne vaut donc $\mathcal{P}_{\text{moy}}=\frac{1}{2}\SI{3}{V}\times\SI{5}{A} \times \frac{1}{2}=\SI{3.75}{W}$.
\end{corrige}

% ************************************** %


% =============================================================================== %
\finEntrainement
% =============================================================================== %





\newpage



% ********************************************************************* %
%                           ENTRAÎNEMENT                                % 
% ********************************************************************* %
% ================ Métadonnées sur l'entraînement ===================== %

\pointInterrogation				{N}
\titreEntrainementFacultatif	{Des calculs de puissance}
\hauteurLargeurCadreReponse		{8mm}{3.5cm}
\dureeResolutionFacultative		{3} % 1, 2, 3 ou 4
\typeEntrainement				{} % Newton ou QCM ou AN ou vide (exercice "normal")
\basiqueEtTransversal			{Y}
\nombreColonnesQuestions		{1} % vide, 1, 2, 3, etc.
\avecPlusieursQuestions			{Y} % Y ou N
\initialisationEntrainement
% ===================================================================== %



On souhaite calculer la puissance reçue par un dipôle. Quand celui-ci est alimenté par une tension $u(t)$ et parcouru par un courant $i(t)$, la puissance moyenne reçue est donnée par la formule :
\[
\mathcal{P}_{\text{moy}}=%\left\langle u(t)\times i(t)\right\rangle =
\dfrac{1}{T}\int_{0}^{T}u(t)\times i(t)\d{}t \minusSmallskip
\]
où $T$ est la période du signal.% et $t_{0}$ un instant quelconque (souvent choisi tel que $t_{0}=0$). 

% =============================================================================== %
\debutEntrainement
% =============================================================================== %

%\medskip
%
%Calculer la puissance moyenne $\mathcal{P}_{\text{moy}}$ du dipôle dans chacun des cas suivants.
%

\bigskip

Dans un premier temps, on considère les signaux
$u(t)=u_{0}\cos\left(\omega t+\psi\right)$ et 
$i(t)=i_{0}\cos\left(\omega t+\psi\right)$.

% ^^^^^^^^^^^^^^^^^^^^^^^^^^^^^^^^^^^^^^ %
%               Question                 %
% ^^^^^^^^^^^^^^^^^^^^^^^^^^^^^^^^^^^^^^ %

\begin{enonce}
Combien vaut la période $T$ pour ces signaux ?
\end{enonce}

\reponse{$\dfrac{2\pi}{\omega}$}

%\begin{corrige}
%\end{corrige}

% ************************************** %



% ^^^^^^^^^^^^^^^^^^^^^^^^^^^^^^^^^^^^^^ %
%               Question                 %
% ^^^^^^^^^^^^^^^^^^^^^^^^^^^^^^^^^^^^^^ %

\begin{enonce}
Calculer $\mathcal{P}_{\text{moy}}$ pour ces signaux. \minusBigskip

{\itshape On pourra utiliser la formule $\cos^2(x)=\dfrac{1+\cos (2x)}{2} $}
\end{enonce}

\reponse{$\dfrac{u_{0}i_{0}}{2}$}

\begin{corrige}
On a $\mathcal{P}(t)=u_{0}i_{0}\cos^{2}\left(\omega t+\psi\right)=\dfrac{u_{0}i_{0}}{2}
\Big(1+\cos\left(2\omega t+2\psi\right)\Big)$.

On intègre:
\begin{align*}
\mathcal{P}_{\text{moy}} 
& =\dfrac{1}{T}\times\dfrac{u_{0}i_{0}}{2}\int_{0}^{T}1+\cos\left(2\omega t+2\psi\right)\d{}t\\
 & =\dfrac{1}{T}\times\dfrac{u_{0}i_{0}}{2}\left[t+\dfrac{1}{2\omega}\sin\left(2\omega t+2\psi\right)\right]_{0}^{T}=\dfrac{u_{0}i_{0}}{2}.
\end{align*}

On peut retenir la propriété $\left\langle \cos^{2}\left(\omega t+\psi\right)\right\rangle =\left\langle \sin^{2}\left(\omega t+\psi\right)\right\rangle =\dfrac{1}{2}$.
\end{corrige}

% ************************************** %



% =============================================================================== %
\pauseEntrainement
% =============================================================================== %

\bigskip\smallskip

Maintenant, on considère les signaux
$u(t)=u_{0}\cos\left(\omega t\right)$
et
$i(t)=i_{0}\cos\left(\omega t+\varphi\right)$.

% =============================================================================== %
\repriseEntrainement
% =============================================================================== %



% ^^^^^^^^^^^^^^^^^^^^^^^^^^^^^^^^^^^^^^ %
%               Question                 %
% ^^^^^^^^^^^^^^^^^^^^^^^^^^^^^^^^^^^^^^ %

\begin{enonce}
Calculer $\mathcal{P}_{\text{moy}}$ pour ces signaux. \minusBigskip

{\itshape On pourra utiliser la formule $\cos (a)\cos (b)=\dfrac{1}{2}\big( \cos (a+b) + \cos (a-b) \big)$}
\end{enonce}

\reponse{$\dfrac{u_{0}i_{0}}{2}\cos(\varphi)$}

\begin{corrige}
On a $\mathcal{P}(t)=u_{0}i_{0}\cos\left(\omega t\right)\cos\left(\omega t+\varphi\right)=\dfrac{u_{0}i_{0}}{2}\left[\cos(\varphi)+\cos\left(2\omega t+\varphi\right)\right]$.

%Comme $\cos\left(2\omega t+\varphi\right)$ est de période $T/2$,
On vérifie ensuite que
\[
\left\langle \cos\left(2\omega t+\varphi\right)\right\rangle 
= \frac{1}{T}\int_0^T \cos\left(2\omega t+\varphi\right) \d t 
= \frac{1}{2\omega T} \Big [ \sin \left(2\omega t+\varphi\right)\Big ]_0^T = 0.
\]
Donc, on a  $\mathcal{P}_{\text{moy}} =\dfrac{u_{0}i_{0}}{2}\cos(\varphi)$.
\end{corrige}

% ************************************** %




% =============================================================================== %
\pauseEntrainement
% =============================================================================== %

\bigskip\smallskip

Enfin, on considère les signaux 
$u(t)=u_{0}\times \big(1+\cos(\omega t)\big)$
et
$i(t)=i_{0}\times \big(2+\sin\left(\omega t+\psi\right)\big)$.

% =============================================================================== %
\repriseEntrainement
% =============================================================================== %




% ^^^^^^^^^^^^^^^^^^^^^^^^^^^^^^^^^^^^^^ %
%               Question                 %
% ^^^^^^^^^^^^^^^^^^^^^^^^^^^^^^^^^^^^^^ %

\begin{enonce}
Calculer $\mathcal{P}_{\text{moy}}$ pour ces signaux.
\end{enonce}

\reponse{$u_{0}i_{0}\left(2+\frac{1}{2}\sin(\psi)\right)$}

\begin{corrigeNewpage}
La puissance peut se décomposer en plusieurs termes:
\begin{align*}
\mathcal{P}(t) & =u_{0}i_{0}\left(1+\cos(\omega t)\right)\left(2+\sin\left(\omega t+\psi\right)\right)\\
 & =u_{0}i_{0}\left(2+2\cos(\omega t)+\sin\left(\omega t+\psi\right)+\cos(\omega t)\sin\left(\omega t+\psi\right)\right)\\
 & =u_{0}i_{0}\left(2+2\cos(\omega t)+\sin\left(\omega t+\psi\right)+\cos(\omega t)\cos\left(\omega t+\psi-\dfrac{\pi}{2}\right)\right)
\end{align*}
On peut alors séparer les calculs de valeurs moyennes:
\begin{align*}
\mathcal{P}_{\text{moy}} & =u_{0}i_{0}\left(2+2\left\langle \cos(\omega t)\right\rangle +\left\langle \sin\left(\omega t+\psi\right)\right\rangle +\left\langle \cos(\omega t)\cos\left(\omega t+\psi-\dfrac{\pi}{2}\right)\right\rangle \right)\\
 & =u_{0}i_{0}\left(2+\frac{1}{2}\cos\left(\psi-\dfrac{\pi}{2}\right)\right)\\
 & =u_{0}i_{0}\left(2+\frac{1}{2}\sin(\psi)\right).
\end{align*}
\end{corrigeNewpage}

% ************************************** %

%================================================================================ %
\finEntrainement
% =============================================================================== %



%
%
%
%% ********************************************************************* %
%%                           ENTRAÎNEMENT                                % 
%% ********************************************************************* %
%% ================ Métadonnées sur l'entraînement ===================== %
%
%\pointInterrogation				{N}
%\titreEntrainementFacultatif	{Calcul de puissance (2)}
%\hauteurLargeurCadreReponse		{8mm}{3.5cm}
%\dureeResolutionFacultative		{2} % 1, 2, 3 ou 4
%\typeEntrainement				{} % Newton ou QCM ou AN ou vide (exercice "normal")
%\nombreColonnesQuestions		{1} % vide, 1, 2, 3, etc.
%\avecPlusieursQuestions			{N} % Y ou N
%\initialisationEntrainement
%% ===================================================================== %
%
%
%Pour déterminer une puissance moyenne, un étudiant a besoin de calculer
%l'intégrale suivante:
%\[
%I=\dfrac{1}{T}\int_{0}^{T}u(t)\cdot i(t)\d{}t
%\]
% où la tension $u(t)$ et l'intensité $i(t)$ prennent respectivement
%les expressions suivantes:
%\[
%u(t)=u_{0}\dfrac{t}{\tau}\exp\left(-\dfrac{t}{\tau}\right)\quad\text{et }\quad i(t)=i_{0}\left(1-\dfrac{t}{\tau}\right)\exp\left(-\dfrac{t}{\tau}\right)
%\]
%
%
%
%% =============================================================================== %
%\debutEntrainement
%% =============================================================================== %
%
%
%%%%%%%%%%%%%%%%%%%%%%%%%%%%%%%%%%%%%%%%%%%%%%%%%%%%%%%%%%
%\begin{enonce}
%S'aidant d'une intégration par parties, établir la puissance moyenne:
%\end{enonce}
%\reponse{$\dfrac{u_{0}i_{0}}{T\tau}\left[\dfrac{T^{2}}{2}\exp\left(-\dfrac{2T}{\tau}\right)\right]=\dfrac{u_{0}i_{0}}{2}T\exp\left(-\dfrac{2T}{\tau}\right)$}
%\begin{corrige}
%La puissance moyenne s'écrit: 
%\begin{align*}
%\mathcal{P}_{\text{moy}} & =\dfrac{1}{T}\int_{0}^{T}u_{0}\left(\dfrac{t}{\tau}\exp\left(-\dfrac{t}{\tau}\right)\right)\times i_{0}\left(\left(1-\dfrac{t}{\tau}\right)\exp\left(-\dfrac{t}{\tau}\right)\right)\d{}t=\dfrac{u_{0}i_{0}}{T\tau}\int_{0}^{T}\left(t-\dfrac{t^{2}}{\tau}\right)\exp\left(-\dfrac{2t}{\tau}\right)\d{}t\\
% & =\dfrac{u_{0}i_{0}}{T\tau}\left[\underset{=I}{\underbrace{\int_{0}^{T}t\exp\left(-\dfrac{2t}{\tau}\right)\d{}t}}-\underset{=J}{\underbrace{\int_{0}^{T}\dfrac{t^{2}}{\tau}\exp\left(-\dfrac{2t}{\tau}\right)\d{}t}}\right]
%\end{align*}
%Pour calculer $I$, on fait l'intégration par partie en posant $\d{}f=t\d{}t$
%et $g=\exp\left(-\dfrac{2t}{\tau}\right)$:
%\begin{align*}
%I & =\left[\dfrac{t^{2}}{2}\exp\left(-\dfrac{2t}{\tau}\right)\right]_{0}^{T}-\intop_{0}^{T}\dfrac{t^{2}}{2}\left(-\dfrac{2}{\tau}\right)\exp\left(-\dfrac{2t}{\tau}\right)\\
% & =\dfrac{T^{2}}{2}\exp\left(-\dfrac{2T}{\tau}\right)+\int_{0}^{T}\dfrac{t^{2}}{\tau}\exp\left(-\dfrac{2t}{\tau}\right)\d{}t\\
% & =\dfrac{T^{2}}{2}\exp\left(-\dfrac{2T}{\tau}\right)+J
%\end{align*}
%Ainsi:
%\[
%\mathcal{P}_{\text{moy}}=\dfrac{u_{0}i_{0}}{T\tau}\left[\dfrac{T^{2}}{2}\exp\left(-\dfrac{2T}{\tau}\right)\right]=\dfrac{u_{0}i_{0}}{2}\dfrac{T}{\tau}\exp\left(-\dfrac{2T}{\tau}\right)
%\]
%
%\end{corrige}
%
%
%%================================================================================ %
%\finEntrainement
%% =============================================================================== %
%
%
%


% ********************************************************************* %
%                           ENTRAÎNEMENT                                % 
% ********************************************************************* %
% ================ Métadonnées sur l'entraînement ===================== %

\pointInterrogation				{N}
\titreEntrainementFacultatif	{Calcul de puissance en RSF}
\hauteurLargeurCadreReponse		{8mm}{3.5cm}
\dureeResolutionFacultative		{2} % 1, 2, 3 ou 4
\typeEntrainement				{} % Newton ou QCM ou AN ou vide (exercice "normal")
\basiqueEtTransversal			{Y}
\nombreColonnesQuestions		{1} % vide, 1, 2, 3, etc.
\avecPlusieursQuestions			{Y} % Y ou N
\initialisationEntrainement
% ===================================================================== %


En régime sinusoïdal forcé, un générateur idéal de tension $\underline{u}$
alimente un dipôle inconnu en délivrant un courant $\text{\ensuremath{\underline{i}}}$. 
Dans ce cas, la puissance moyenne peut être calculée à l'aide de la formule
\[
\mathcal{P}_{\text{moy}}=\dfrac{1}{2}\operatorname{Re}\left(\underline{u}\cdot \underline{i}^{\star}\right)=\dfrac{1}{2}\operatorname{Re}\left(\underline{u}^{\star}\cdot \underline{i}\right)
\]
 où $\underline{x}^{\star}$ est le complexe conjugué de $\underline{x}$.



% =============================================================================== %
\debutEntrainement
% =============================================================================== %

\smallskip

Exprimer la puissance moyenne reçue par le dipôle quand:

% ************************************** %



% ^^^^^^^^^^^^^^^^^^^^^^^^^^^^^^^^^^^^^^ %
%               Question                 %
% ^^^^^^^^^^^^^^^^^^^^^^^^^^^^^^^^^^^^^^ %

\begin{enonce}
$\underline{u}=U\eEuler^{\jC \omega t}$ \quad et\quad $\underline{i}=jC\omega\underline{u}$
\end{enonce}

\reponse{0}

\begin{corrige}
On a $\mathcal{P}_{\text{moy}}=\dfrac{1}{2}\operatorname{Re}\left(jC\omega\left|\underline{u}\right|^{2}\right)=0$.
\end{corrige}

% ************************************** %



% ^^^^^^^^^^^^^^^^^^^^^^^^^^^^^^^^^^^^^^ %
%               Question                 %
% ^^^^^^^^^^^^^^^^^^^^^^^^^^^^^^^^^^^^^^ %

\begin{enonce}
$\underline{i}=I\eEuler^{\jC \omega t+\varphi}$  \quad et\quad $\underline{u}=jL\omega\underline{i}$
\end{enonce}

\reponse{0}

\begin{corrige}
On a $\mathcal{P}_{\text{moy}}=\dfrac{1}{2}\operatorname{Re}\left(jL\omega\left|\underline{i}\right|^{2}\right)=0$.
\end{corrige}

% ************************************** %



% ^^^^^^^^^^^^^^^^^^^^^^^^^^^^^^^^^^^^^^ %
%               Question                 %
% ^^^^^^^^^^^^^^^^^^^^^^^^^^^^^^^^^^^^^^ %

\begin{enonce}
$\underline{u}=\sqrt{2}\left(1-j\right)\eEuler^{\jC \omega t}$  \quad et\quad $\underline{i}=3\left(\dfrac{1}{2}+j\dfrac{\sqrt{3}}{2}\right)\eEuler^{\jC \omega t}$
\end{enonce}

\reponse{$3\cos\left(\dfrac{7\pi}{12}\right)\,\si{W}$}

\begin{corrige}
Commençons par réécrire $\underline{u}$ et $\underline{i}$:
\begin{align*}
\underline{u} & =2\left(\dfrac{1}{\sqrt{2}}-\dfrac{j}{\sqrt{2}}\right)\eEuler^{\jC \omega t}=2\left(\eEuler^{-j\dfrac{\pi}{4}}\right)\eEuler^{\jC \omega t}=2\eEuler^{\jC \left(\omega t-\pi/4\right)}\\
\underline{i} & =3\left(\eEuler^{\jC \dfrac{\pi}{3}}\right)\eEuler^{\jC \omega t}=3\eEuler^{\jC \left(\omega t+\pi/3\right)}.
\end{align*}
On en déduit
$\mathcal{P}_{\text{moy}}=\dfrac{1}{2}\operatorname{Re}\left(6\eEuler^{\jC \left(\omega t-\pi/4\right)}\times \eEuler^{\jC \left(\omega t+\pi/3\right)}\right)=3\operatorname{Re}\left(\eEuler^{-j\left(\pi/3+\pi/4\right)}\right)=3\cos\left(\dfrac{7\pi}{12}\right)\,\si{W}$.
\end{corrige}

% ************************************** %



% ^^^^^^^^^^^^^^^^^^^^^^^^^^^^^^^^^^^^^^ %
%               Question                 %
% ^^^^^^^^^^^^^^^^^^^^^^^^^^^^^^^^^^^^^^ %

\begin{enonce}
$\underline{u}=4\sqrt{2}\eEuler^{\jC \left(\omega t+\pi/4\right)}$  \quad et\quad $\underline{i}=\left(3+5j\right)\eEuler^{\jC \omega t}$
\end{enonce}

\reponse{$\SI{16}{W}$}

\begin{corrige}
On a 
\begin{align*}
\mathcal{P}_{\text{moy}} 
& =\dfrac{1}{2}\operatorname{Re}\left(4\sqrt{2}\eEuler^{\jC \left(\omega t+\pi/4\right)}\times\left(3-5j\right)\eEuler^{-j\omega t}\right)=2\sqrt{2}\operatorname{Re}\left(\left(3-5j\right)\eEuler^{\jC \pi/4}\right)\\
 & =2\sqrt{2}\left(\dfrac{3}{\sqrt{2}}+\dfrac{5}{\sqrt{2}}\right)=\SI{16}{W}.
\end{align*}
\end{corrige}

% ************************************** %


%================================================================================ %
\finEntrainement
% =============================================================================== %




\newpage



% •••••••••••••••••••••••••••••••••••••••••••••••••••••••••••••••••••••••••••• %
% •••••••••••••••••••••••••••••••••••••••••••••••••••••••••••••••••••••••••••• %
\sectionFicheEntrainement{Régime permanent}
% •••••••••••••••••••••••••••••••••••••••••••••••••••••••••••••••••••••••••••• %
% •••••••••••••••••••••••••••••••••••••••••••••••••••••••••••••••••••••••••••• %


% ********************************************************************* %
%                           ENTRAÎNEMENT                                % 
% ********************************************************************* %
% ================ Métadonnées sur l'entraînement ===================== %

\titreEntrainementFacultatif	{Puissance consommée}
\hauteurLargeurCadreReponse		{8mm}{4cm}
\dureeResolutionFacultative		{1} % 1, 2, 3 ou 4
\typeEntrainement				{} % Newton ou QCM ou AN ou vide (exercice "normal")
\basiqueEtTransversal			{Y}
\nombreColonnesQuestions		{1} % vide, 1, 2, 3, etc.
\avecPlusieursQuestions			{Y} % Y ou N
\initialisationEntrainement
% ===================================================================== %


% mmmmmmmmmmmmmmmmmmmmmmmmmmmmmmmmmmmmmmmmmmmmmmmmmmmmmmmmm %
% 		  Insertion d'une image en regard d'un texte
% mmmmmmmmmmmmmmmmmmmmmmmmmmmmmmmmmmmmmmmmmmmmmmmmmmmmmmmmm %
% --------------------------------------------------------- %
\pourcentageDeLaPartieAGauche   {0.55}
% --------------------------------------------------------- %
%                                                           %
% -------------------- Partie à gauche -------------------- %
%                                                           %
                                \initialisationPartieGauche % 
%                                                           %
%                                                           %

Soit un générateur réel de fém $E$ constante et de résistance interne $r$. 

On branche à ses bornes un conducteur ohmique de résistance variable $R$.


% --------------------------------------------------------- %
%                                                           %
% -------------------- Partie à droite -------------------- %
%                                                           %
                                \initialisationPartieDroite %
%                                                           %
%                                                           %
\begin{center}
\subimport{_images/}{fiche_ELC04_adaptation_impedance_continu}
\end{center}
% --------------------------------------------------------- %
                               \finalisationDuPartageDePage %					
% --------------------------------------------------------- %


\smallskip

% =============================================================================== %
\debutEntrainement
% =============================================================================== %


% ^^^^^^^^^^^^^^^^^^^^^^^^^^^^^^^^^^^^^^ %
%               Question                 %
% ^^^^^^^^^^^^^^^^^^^^^^^^^^^^^^^^^^^^^^ %

\begin{enonce}
Déterminer l'intensité du courant qui circule dans le circuit
\end{enonce}

\reponse{$\frac{E}{r+R}$}

\begin{corrige}
La loi des mailles permet d'écrire $E=u_r+u_R=rI+RI=(r+R)I$. On a donc $I=\frac{E}{r+R}$.
\end{corrige}

% ************************************** %


% ^^^^^^^^^^^^^^^^^^^^^^^^^^^^^^^^^^^^^^ %
%               Question                 %
% ^^^^^^^^^^^^^^^^^^^^^^^^^^^^^^^^^^^^^^ %

\begin{enonce}
Déterminer la puissance $\mathcal{P}$ dissipée dans le conducteur ohmique en fonction de $E$, $r$ et $R$.

\end{enonce}

\reponse{$E^2\frac{R}{(r+R)^2}$}

\begin{corrige}
La puissance dissipée dans le conducteur ohmique de résistance $R$ vaut $\mathcal{P}=u_R I=RI^2=E^2\frac{R}{(r+R)^2}$.
\end{corrige}

% ************************************** %

% =============================================================================== %
\finEntrainement
% =============================================================================== %





% ********************************************************************* %
%                           ENTRAÎNEMENT                                % 
% ********************************************************************* %
% ================ Métadonnées sur l'entraînement ===================== %

\titreEntrainementFacultatif	{Optimisation de puissance échangée}
\hauteurLargeurCadreReponse		{8mm}{4cm}
\dureeResolutionFacultative		{1} % 1, 2, 3 ou 4
\typeEntrainement				{} % Newton ou QCM ou AN ou vide (exercice "normal")
\basiqueEtTransversal			{Y}
\nombreColonnesQuestions		{1} % vide, 1, 2, 3, etc.
\avecPlusieursQuestions			{Y} % Y ou N
\initialisationEntrainement
% ===================================================================== %


Dans un certain circuit, la puissance dissipée dans un conducteur ohmique de résistance $R$ %du circuit de l'entraînement  précédent 
vaut 
\[
\mathcal{P}=E^2\frac{R}{(r+R)^2},
\]
où $r$ est un paramètre.

On souhaite déterminer quelle valeur de $R$ permet d'optimiser la puissance reçue par la résistance $R$ étant données les caractéristiques de la source 

% =============================================================================== %
\debutEntrainement
% =============================================================================== %




%^^^^^^^^^^^^^^^^^^^^^^ %

\begin{enonce}
Calculer $\diff{\mathcal{P}}{R}$
\end{enonce}

\reponse{$E^2\frac{r-R}{\left(r+R \right)^3}$}

\begin{corrige}
Il faut dériver la fonction $\mathcal{P}(R)$. On calcule% qui est un quotient. On utilise la formule de dérivation du quotient $\left( \frac{u}{v} \right)'=\frac{u'v-uv'}{v^2}$ avec ici $u=R$, $u'=1$, $v=(r+R)^2$ et $v'=2(r+R)$, soit
\[
\diff{\mathcal{P}}{R}=E^2\frac{1 \times (r+R)^2-R\times 2(r+R)}{\left(r+R \right)^4}=E^2(r+R)\frac{(r+R)-2R}{\left(r+R \right)^4}
\]
soit finalement
\[
\diff{\mathcal{P}}{R}=E^2(r+R)\frac{r-R}{\left(r+R \right)^4}=E^2\frac{r-R}{\left(r+R \right)^3}.
\]
\end{corrige}

% ************************************** %


% ^^^^^^^^^^^^^^^^^^^^^^^^^^^^^^^^^^^^^^ %
%               Question                 %
% ^^^^^^^^^^^^^^^^^^^^^^^^^^^^^^^^^^^^^^ %

\begin{enonce}
Trouver la valeur $R_\text{max}$ pour laquelle $\mathcal{P}(R)$ est maximale
\begin{listeQCM4Colonnes}
	\item $R_\text{max}=R$ \smallskip
	\item $R_\text{max}=r$
	\columnbreak
	\item $R_\text{max}=R+r$
	\item $R_\text{max}=\frac{R^2}{r+R}$
	\columnbreak
	\item $R_\text{max}=\frac{1}{\frac{1}{r}+\frac{1}{R}}$
	\columnbreak
	\item $R_\text{max}=R\times \eEuler^{r/R}$
	\smallskip
\end{listeQCM4Colonnes}

\bigskip

\end{enonce}

\reponse{\reponseB{}}

\begin{corrigeNewpage}
Il faut annuler la dérivée pour trouver l'extremum de $\mathcal{P}(R)$. Comme $\mathcal{P}(R)$ est positive et vaut $0$ en $R=0$ et en $R\rightarrow \infty$, alors cet extremum est un maximum. 
On a alors, par annulation du numérateur $R_\text{max}=r$.
\end{corrigeNewpage}

% ************************************** %


% =============================================================================== %
\finEntrainement
% =============================================================================== %





% ********************************************************************* %
%                           ENTRAÎNEMENT                                % 
% ********************************************************************* %
% ================ Métadonnées sur l'entraînement ===================== %

\titreEntrainementFacultatif	{Un peu de calcul algébrique}
\hauteurLargeurCadreReponse		{8mm}{4cm}
\dureeResolutionFacultative		{1} % 1, 2, 3 ou 4
\typeEntrainement				{} % Newton ou QCM ou AN ou vide (exercice "normal")
\basiqueEtTransversal			{Y}
\nombreColonnesQuestions		{1} % vide, 1, 2, 3, etc.
\avecPlusieursQuestions			{N} % Y ou N
\initialisationEntrainement
% ===================================================================== %


On considère une résistance $R$ définie par
\[
R=R_0\times \eEuler^{r/R_0}.
\]



% =============================================================================== %
\debutEntrainement
% =============================================================================== %


% ^^^^^^^^^^^^^^^^^^^^^^^^^^^^^^^^^^^^^^ %
%               Question                 %
% ^^^^^^^^^^^^^^^^^^^^^^^^^^^^^^^^^^^^^^ %

\begin{enonce}
Déterminer, en fonction de $R_0$, la valeur du paramètre $r$ pour que $R = 2R_0$.
\bigskip

\end{enonce}

\reponse{$\ln(2) R_0$}

\begin{corrige}
Si on a $R = 2R_0$, alors on a $\eEuler^{r/R_0} =2$ et donc $r/R_0 = \ln(2)$. Finalement, on a 
$r = \ln(2) R_0$.
\end{corrige}

% ************************************** %


% =============================================================================== %
\finEntrainement
% =============================================================================== %






\newpage



% ********************************************************************* %
%                           ENTRAÎNEMENT                                % 
% ********************************************************************* %
% ================ Métadonnées sur l'entraînement ===================== %

\titreEntrainementFacultatif	{Charge d'une batterie}
\hauteurLargeurCadreReponse		{8mm}{4cm}
\dureeResolutionFacultative		{3} % 1, 2, 3 ou 4
\typeEntrainement				{} % Newton ou QCM ou AN ou vide (exercice "normal")
\nombreColonnesQuestions		{1} % vide, 1, 2, 3, etc.
\avecPlusieursQuestions			{Y} % Y ou N
\initialisationEntrainement
% ===================================================================== %


Une batterie de voiture est déchargée. 
Pour recharger cette batterie, de fém $e=\SI{12}{\volt}$ et de résistance interne $r=\SI{0,2}{\ohm}$, on la branche sur un chargeur de fém $E=\SI{13}{\volt}$ et de résistance interne $R=\SI{0,3}{\ohm}$. 

On a alors le circuit suivant : 


\begin{center}
\subimport{_images/}{fiche_ELC04_charge_batterie}
\end{center}

On lit sur la batterie qu'elle a une capacité de $\SI{50}{\ampere \hour}$ (ampères-heures). 

\bigskip



% =============================================================================== %
\debutEntrainement
% =============================================================================== %

% ^^^^^^^^^^^^^^^^^^^^^^^^^^^^^^^^^^^^^^ %
%               Question                 %
% ^^^^^^^^^^^^^^^^^^^^^^^^^^^^^^^^^^^^^^ %

\begin{enonce}
Exprimer le courant $I$ circulant dans la batterie en fonction de $E$, $e$, $R$ et $r$.

\end{enonce}

\reponse{$\frac{E-e}{R+r}$}

\begin{corrige}
On applique la loi des mailles $E-U_R-u_r-e=0$. On a donc $E-e=(R+r)I$, et donc $I=\frac{E-e}{R+r}$.
\end{corrige}

% ************************************** %


% ^^^^^^^^^^^^^^^^^^^^^^^^^^^^^^^^^^^^^^ %
%               Question                 %
% ^^^^^^^^^^^^^^^^^^^^^^^^^^^^^^^^^^^^^^ %

\begin{enonce}
Exprimer la tension $U$ aux bornes de la batterie lors de la charge  en fonction de $E$, $e$, $R$ et $r$.

\end{enonce}

\reponse{$\frac{eR+Er}{R+r}$}

\begin{corrige}
La batterie est en convention récepteur donc on a
\[
U=e+rI=e+r\frac{E-e}{R+r}=\frac{eR+er+rE-re}{R+r}=\frac{eR+Er}{R+r}.
\]
\end{corrige}

% ************************************** %


% ^^^^^^^^^^^^^^^^^^^^^^^^^^^^^^^^^^^^^^ %
%               Question                 %
% ^^^^^^^^^^^^^^^^^^^^^^^^^^^^^^^^^^^^^^ %

\begin{enonce}
Exprimer la puissance délivrée par la source de fém $E$ en fonction de $E$, $e$, $R$ et $r$.

\end{enonce}

\reponse{$E\frac{E-e}{R+r}$}

\begin{corrige}
La puissance fournie par le chargeur vaut $\mathcal{P}=EI=E\frac{E-e}{R+r}$.
\end{corrige}

% ************************************** %


% ^^^^^^^^^^^^^^^^^^^^^^^^^^^^^^^^^^^^^^ %
%               Question                 %
% ^^^^^^^^^^^^^^^^^^^^^^^^^^^^^^^^^^^^^^ %

\begin{enonce}
Exprimer la puissance dissipée par effet Joule dans le circuit  en fonction de $E$, $e$, $R$ et $r$.

\end{enonce}

\reponse{$\frac{(E-e)^2}{R+r}$}

\begin{corrige}
La puissance est dissipée par effet Joule dans les deux conducteurs ohmiques et vaut donc 
\[
\mathcal{P}_J=RI^2+rI^2=(R+r)\left(\frac{E-e}{R+r}\right)^2 =\frac{(E-e)^2}{R+r}.
\]
\end{corrige}

% ************************************** %


% ^^^^^^^^^^^^^^^^^^^^^^^^^^^^^^^^^^^^^^ %
%               Question                 %
% ^^^^^^^^^^^^^^^^^^^^^^^^^^^^^^^^^^^^^^ %

\begin{enonce}
Exprimer la puissance reçue par la batterie en fonction de $E$, $e$, $R$ et $r$.

\end{enonce}

\reponse{$e\frac{E-e}{R+r}$}

\begin{corrige}
La puissance reçue par la batterie vaut $\mathcal{P}=eI=e\frac{E-e}{R+r}$ car elle est en convention récepteur.
\end{corrige}

% ************************************** %




% =============================================================================== %
\debutEntrainement
% =============================================================================== %

\bigskip

Le rendement $\eta$ de la charge est égal au rapport de la puissance reçue par la batterie par la puissance fournie par la source $E$.


% =============================================================================== %
\finEntrainement
% =============================================================================== %




% ^^^^^^^^^^^^^^^^^^^^^^^^^^^^^^^^^^^^^^ %
%               Question                 %
% ^^^^^^^^^^^^^^^^^^^^^^^^^^^^^^^^^^^^^^ %

\begin{enonce}
Déterminer l'expression du rendement $\eta$  en fonction de $E$ et  $e$.

\end{enonce}

\reponse{$\frac{e}{E}$}

\begin{corrige}
En suivant la définition de l'énoncé, on trouve
\[
\eta=\frac{e\frac{E-e}{R+r}}{E\frac{E-e}{R+r}}=\frac{e}{E}.
\]
\end{corrige}

% ************************************** %


% ^^^^^^^^^^^^^^^^^^^^^^^^^^^^^^^^^^^^^^ %
%               Question                 %
% ^^^^^^^^^^^^^^^^^^^^^^^^^^^^^^^^^^^^^^ %

\begin{enonce}
Calculer la valeur numérique du rendement $\eta$
\end{enonce}

\reponse{$92 \, \% $}

\begin{corrige}
Numériquement, on calcule $\eta=12/13=92 \, \%$.
\end{corrige}

% ************************************** %


\newpage



% ********************************************************************* %
%                           ENTRAÎNEMENT                                % 
% ********************************************************************* %
% ================ Métadonnées sur l'entraînement ===================== %

\pointInterrogation				{N}
\titreEntrainementFacultatif	{Énergie d'un condensateur en régime permanent}
\hauteurLargeurCadreReponse		{8mm}{3.5cm}
\dureeResolutionFacultative		{2} % 1, 2, 3 ou 4
\typeEntrainement				{QCM} % Newton ou QCM ou AN ou vide (exercice "normal")
\nombreColonnesQuestions		{1} % vide, 1, 2, 3, etc.
\avecPlusieursQuestions			{Y} % Y ou N
\initialisationEntrainement
% ===================================================================== %

En régime permanent, l'énergie stockée dans un condensateur de capacité $C$ est $\mathcal{E}=\dfrac{1}{2}Cu^{2}$, où $u$ est  la tension à ses bornes.

\minusSmallskip

% =============================================================================== %
\debutEntrainement
% =============================================================================== %

\begin{center}
\subimport{_images/}{fiche_ELC02_Condo_energie_E_CBD_v2.tex}
\end{center}


\minusMedskip

% ^^^^^^^^^^^^^^^^^^^^^^^^^^^^^^^^^^^^^^ %
%               Question                 %
% ^^^^^^^^^^^^^^^^^^^^^^^^^^^^^^^^^^^^^^ %

\begin{enonce}
On considère que le régime permanent est établi dans le circuit \no{}1.

Dans quel condensateur l'énergie stockée est-elle la plus importante?

\begin{listeQCM4Colonnes}
\item $C$
\item $2C$
\item $3C$
\end{listeQCM4Colonnes}
\minusMedskip
\end{enonce}

\reponse{\reponseC{}}

\begin{corrige}
On fait le schéma
\begin{center}
\shorthandoff{:!} %règle un problème de compatibilité avec babel
%\begin{circuitikz}[straight voltages,scale=0.9]
\begin{circuitikz}[scale=0.7]
\begin{scope}
\draw (-1,0) to[vsource,V=$E$] (-1,3.5) -- (1.5,3.5) to[C,l=$C$,v<=$U_1$] (1.5,0) -- (-1,0);
\draw (1.5,3.5)  to[R,l=$R$] (4,3.5) to[R,l=$2R$] (4,0);
\draw (6.5,3.5)  to[C,l=$2C$,v<=$U_2$] (6.5,0);
\draw (8,3.5)  to[C,l=$3C$,v<=$U_3$] (8,1.75) to[R,l=$R$] (8,0) -- (1.5,0);
\draw (4,3.5) -- (8,3.5);
\end{scope}

\end{circuitikz}
\shorthandon{:!}
\end{center}
En régime permanent, les condensateurs se comportent comme des interrupteurs
ouverts:
\[
U_{1}=E\quad U_{2}=U_{3}=\dfrac{2}{3}E
\]
les énergies stockées dans les condensateurs sont alors:
\[
\mathcal{E}_{1}=\dfrac{1}{2}CE^{2}\quad\mathcal{E}_{2}=\dfrac{4}{9}CE^{2}\quad\mathcal{E}_{3}=\dfrac{3}{2}C\left(\dfrac{4}{9}E^{2}\right)=\dfrac{2}{3}CE^{2}
\]
on a alors $\mathcal{E}_{2}<\mathcal{E}_{1}<\mathcal{E}_{3}$. C'est la réponse \reponseC{} qui est la bonne.
\end{corrige}

% ************************************** %


\minusSmallskip

% ^^^^^^^^^^^^^^^^^^^^^^^^^^^^^^^^^^^^^^ %
%               Question                 %
% ^^^^^^^^^^^^^^^^^^^^^^^^^^^^^^^^^^^^^^ %

\begin{enonce}
Même question pour le circuit \no{}2.
\begin{listeQCM4Colonnes}
\item $C$
\item $2C$
\item $3C$
\end{listeQCM4Colonnes}
\minusMedskip

\end{enonce}
\reponse{\reponseA{}}

\begin{corrige}
On fait le schéma
\begin{center}
%\subimport{_images/}{fiche_ELC02_Condo_energie_C2.tex}
\shorthandoff{:!} %règle un problème de compatibilité avec babel
\begin{circuitikz}[scale=0.7]
\begin{scope}
\draw (-1,0) to[vsource,V=$\SI{6}{V}$] (-1,3) -- (1.5,3) to[C,l=$3C$,v<=$U_3$] (1.5,0) -- (-1,0);
\draw (1.5,3) to[C,l=$2C$,v<=$U_2$] (3.5,3) to[R,l=$R$] (5.5,3);
\draw (5.5,0) to[vsource,V=$\SI{12}{V}$] (5.5,3) -- (8,3) to[C,l=$C$,v<=$U_1$] (8,0) -- (1.5,0);
\end{scope}

\end{circuitikz}
\shorthandon{:!}
\end{center}

En régime permanent, les condensateurs se comportent comme des interrupteurs
ouverts:
\[
U_{1}=\SI{12}{V}\quad U_{2}=\SI{-6}{V}\quad U_{3}=\SI{6}{V}
\]
les énergies stockées dans les condensateurs est alors:
\[
\mathcal{E}_{1}=\dfrac{1}{2}C\left(12\right)^{2}=72C\quad\mathcal{E}_{2}=\dfrac{1}{2}\times2C\left(6\right)^{2}=36C\quad\mathcal{E}_{3}=\dfrac{1}{2}\times3C\left(6\right)^{2}=54C
\]
on a alors $\mathcal{E}_{1}>\mathcal{E}_{3}>\mathcal{E}_{2}$. C'est la réponse \reponseA{} qui est la bonne.
\end{corrige}

% ************************************** %



%================================================================================ %
\finEntrainement
% =============================================================================== %









% ********************************************************************* %
%                           ENTRAÎNEMENT                                % 
% ********************************************************************* %
% ================ Métadonnées sur l'entraînement ===================== %

\pointInterrogation				{N}
\titreEntrainementFacultatif	{Énergie d'une bobine en régime permanent}
\hauteurLargeurCadreReponse		{8mm}{3.5cm}
\dureeResolutionFacultative		{3} % 1, 2, 3 ou 4
\typeEntrainement				{QCM} % Newton ou QCM ou AN ou vide (exercice "normal")
\nombreColonnesQuestions		{1} % vide, 1, 2, 3, etc.
\avecPlusieursQuestions			{Y} % Y ou N
\initialisationEntrainement
% ===================================================================== %

En régime permanent, l'énergie stockée dans une bobine d'inductance $L$ est $\mathcal{E}=\dfrac{1}{2}Li^{2}$ où $i$ est le courant qui la traverse.

\minusBigskip\minusSmallskip


% =============================================================================== %
\debutEntrainement
% =============================================================================== %

\begin{center}
\subimport{_images/}{fiche_ELC02_Bob_energie_E_CBD_v2.tex}
\end{center}

\minusSmallskip\minusSmallskip

% ^^^^^^^^^^^^^^^^^^^^^^^^^^^^^^^^^^^^^^ %
%               Question                 %
% ^^^^^^^^^^^^^^^^^^^^^^^^^^^^^^^^^^^^^^ %

\begin{enonce}
On considère que le régime permanent est établi dans le circuit \no 3.

Dans quelle bobine l'énergie stockée est-elle la plus importante?

\begin{listeQCM4Colonnes}
\item $L$
\item $2L$
\item $3L$
\end{listeQCM4Colonnes}
\minusMedskip

\end{enonce}

\reponse{\reponseB{}}

\begin{corrige}
L'énergie stockée dans les différentes bobines est:
\[
\mathcal{E}_{1}=\dfrac{1}{2}L\left(4\right)^{2}=8L\quad\mathcal{E}_{2}=\dfrac{1}{2}\times2L\left(3\right)^{2}=9L\quad\mathcal{E}_{3}=\dfrac{1}{2}\times3L\left(1\right)^{2}=\dfrac{3}{2}L
\]
 donc $\mathcal{E}_{3}<\mathcal{E}_{1}<\mathcal{E}_{2}$. C'est la réponse \reponseB{} qui est la bonne.
\end{corrige}

% ************************************** %


\minusSmallskip


% ^^^^^^^^^^^^^^^^^^^^^^^^^^^^^^^^^^^^^^ %
%               Question                 %
% ^^^^^^^^^^^^^^^^^^^^^^^^^^^^^^^^^^^^^^ %

\begin{enonce}
Même question pour le circuit \no 4.
\begin{listeQCM4Colonnes}
\item $L$
\item $2L$
\item $3L$
\end{listeQCM4Colonnes}
\minusMedskip

\end{enonce}

\reponse{\reponseA{}}

\begin{corrige}
Les bobines se comportent comme des fils en régime permanent. Le montage
se simplifie alors:

\begin{center}
%\subimport{_images/}{fiche_ELC02_Bob_energie_C.tex}
\shorthandoff{:!} %règle un problème de compatibilité avec babel
\begin{circuitikz}[scale=0.8]
\draw (0,0) to[vsource,V=$E$] (0,2.5)  to[short,i=$I_1$] (2,2.5) to[R,l=$R$] (2,0);
\draw (2,2.5) to[R,l=$R$] (4,2.5)  to[R,l=$R$,i=$I_2$] (4,0);
\draw (5.5,2.5)  to[R,l=$2R$,i=$I_3$] (5.5,0);
\draw (4,2.5) -- (5.5,2.5);
\draw (5.5,0) -- (0,0);
\end{circuitikz}
\shorthandon{:!}
\end{center}

En calculant les résistances équivalentes, on peut déterminer les valeurs des courants:
\[
I_{1}=\dfrac{8E}{5R},\qquad I_{2}=\dfrac{2}{3}\left(\dfrac{3E}{5R}\right)=\dfrac{2E}{5R}\qquad\text{et}\qquad I_{3}=\dfrac{1}{3}\left(\dfrac{3E}{5R}\right)=\dfrac{E}{5R}.
\]
Ainsi les énergies stockées dans les bobines sont:
\[
\mathcal{E}_{1}=\dfrac{1}{2}L\left(\dfrac{8E}{5R}\right)^{2}=\dfrac{32}{25}\dfrac{LE^{2}}{R^{2}}\qquad\mathcal{E}_{2}=\dfrac{1}{2}\times2L\left(\dfrac{2E}{5R}\right)^{2}=\dfrac{4}{25}\dfrac{LE^{2}}{R^{2}}\qquad\text{et}\qquad\mathcal{E}_{3}=\dfrac{1}{2}\times3L\left(\dfrac{E}{5R}\right)^{2}=\dfrac{3}{50}\dfrac{LE^{2}}{R^{2}}.
\]
On a $\mathcal{E}_{3}<\mathcal{E}_{2}<\mathcal{E}_{1}$ : c'est la réponse \reponseA{} qui est la bonne.
\end{corrige}

% ************************************** %



%================================================================================ %
\finEntrainement
% =============================================================================== %








\newpage




% •••••••••••••••••••••••••••••••••••••••••••••••••••••••••••••••••••••••••••• %
% •••••••••••••••••••••••••••••••••••••••••••••••••••••••••••••••••••••••••••• %
\sectionFicheEntrainement{Bilans d'énergie pour des circuits soumis à des échelons de tension}
% •••••••••••••••••••••••••••••••••••••••••••••••••••••••••••••••••••••••••••• %
% •••••••••••••••••••••••••••••••••••••••••••••••••••••••••••••••••••••••••••• %

\medskip

\begin{prerequis}
L'énergie $\mathcal{E}$ fournie à un dipôle entre les temps $t_0$ et $t_1$ est égale à
\[
\mathcal{E}
=
\int_{t_0}^{t_1} \mathcal{P}(t)\d t
\]
où $\mathcal{P}(t)$ est la puissance  instantanée fournie à ce dipôle.
\end{prerequis}




% ********************************************************************* %
%                           ENTRAÎNEMENT                                % 
% ********************************************************************* %
% ================ Métadonnées sur l'entraînement ===================== %

\titreEntrainementFacultatif	{Charge d'un condensateur}
\hauteurLargeurCadreReponse		{8mm}{4cm}
\dureeResolutionFacultative		{3} % 1, 2, 3 ou 4
\typeEntrainement				{} % Newton ou QCM ou AN ou vide (exercice "normal")
\nombreColonnesQuestions		{1} % vide, 1, 2, 3, etc.
\avecPlusieursQuestions			{Y} % Y ou N
\initialisationEntrainement
% ===================================================================== %



% mmmmmmmmmmmmmmmmmmmmmmmmmmmmmmmmmmmmmmmmmmmmmmmmmmmmmmmmm %
% 		  Insertion d'une image en regard d'un texte
% mmmmmmmmmmmmmmmmmmmmmmmmmmmmmmmmmmmmmmmmmmmmmmmmmmmmmmmmm %
% --------------------------------------------------------- %
\pourcentageDeLaPartieAGauche   {0.55}
% --------------------------------------------------------- %
%                                                           %
% -------------------- Partie à gauche -------------------- %
%                                                           %
                                \initialisationPartieGauche % 
%                                                           %
%                                                           %
Soit le circuit ci-contre dans lequel le condensateur $C$ est initialement déchargé. 

À $t=0$, on ferme l'interrupteur $K$.

Dans ces conditions, la tension aux bornes du condensateur vaut 
\[
u_C(t)=E(1-\exp(-t/\tau))
\]
 avec $\tau=RC$ ; l'intensité dans le circuit vaut 
 \[
 i(t)=\frac{CE}{\tau}\exp(-t/\tau).
 \]


% --------------------------------------------------------- %
%                                                           %
% -------------------- Partie à droite -------------------- %
%                                                           %
                                \initialisationPartieDroite %
%                                                           %
%                                                           %
\begin{center}
\subimport{_images/}{fiche_ELC04_charge_condensateur}
\end{center}
% --------------------------------------------------------- %
                               \finalisationDuPartageDePage %					
% --------------------------------------------------------- %



% =============================================================================== %
\debutEntrainement
% =============================================================================== %

\bigskip

Exprimer, en fonction des grandeurs introduites :

% ^^^^^^^^^^^^^^^^^^^^^^^^^^^^^^^^^^^^^^ %
%               Question                 %
% ^^^^^^^^^^^^^^^^^^^^^^^^^^^^^^^^^^^^^^ %

\begin{enonce}
la puissance instantanée $\mathcal{P}_E(t)$ délivrée par la source de fém $E$.
\end{enonce}

\reponse{$\frac{CE^2}{\tau}\exp(-t/\tau)$}

\begin{corrige}
La puissance instantanée délivrée par la source vaut 
\[
\mathcal{P}_E(t)=Ei(t)=E\times \dfrac{CE}{\tau}\exp(-t/\tau)=\frac{CE^2}{\tau}\exp(-t/\tau).
\]
\end{corrige}

% ************************************** %


% ^^^^^^^^^^^^^^^^^^^^^^^^^^^^^^^^^^^^^^ %
%               Question                 %
% ^^^^^^^^^^^^^^^^^^^^^^^^^^^^^^^^^^^^^^ %

\begin{enonce}
la puissance instantanée $\mathcal{P}_J(t)$ dissipée par effet Joule dans le circuit.
\end{enonce}

\reponse{$\frac{CE^2}{\tau}\exp(-2t/\tau)$}

\begin{corrige}
La puissance dissipée par effet Joule l'est dans la résistance et vaut donc 
\[
\mathcal{P}_J(t)=Ri^2(t)=\frac{R(CE)^2}{\tau^2}\exp(-2t/\tau).
\]
En simplifiant à l'aide de la relation $\tau=RC$, on trouve $\mathcal{P}_J(t)=\frac{CE^2}{\tau}\exp(-2t/\tau)$.
\end{corrige}

% ************************************** %


% ^^^^^^^^^^^^^^^^^^^^^^^^^^^^^^^^^^^^^^ %
%               Question                 %
% ^^^^^^^^^^^^^^^^^^^^^^^^^^^^^^^^^^^^^^ %

\begin{enonce}
la puissance instantanée $\mathcal{P}_C(t)$ reçue par le condensateur
\end{enonce}

\reponse{$\frac{CE^2}{\tau}\Big(\exp(-t/\tau)-\exp(-2t/\tau)\Big)$}

\begin{corrige}
La puissance instantanée reçue par le condensateur vaut 
\[
\mathcal{P}_C(t)
= u_C(t)i(t)=E(1-\exp(-t/\tau)) \times \frac{CE}{\tau}\exp(-t/\tau)
= \frac{CE^2}{\tau}\Big(\exp(-t/\tau)-\exp(-2t/\tau)\Big).
\]
Remarquons que par conservation de la puissance, cette dernière expression peut s'obtenir en faisant la différence entre les deux précédentes, la puissance reçue par le condensateur étant égale à la puissance fournie par la source de tension dont on a retranché la puissance dissipée dans le conducteur ohmique. C'est un bon moyen de contrôler le résultat.
\end{corrige}

% ************************************** %



% ^^^^^^^^^^^^^^^^^^^^^^^^^^^^^^^^^^^^^^ %
%               Question                 %
% ^^^^^^^^^^^^^^^^^^^^^^^^^^^^^^^^^^^^^^ %

\begin{enonce}
l'énergie totale $\mathcal{E}_E$ fournie par la source de tension que l'on calculera grâce à la formule \minusSmallskip
\[
\mathcal{E}_E = \int_0^\infty  \mathcal{P}_E(t) \d t.\minusBigskip\minusSmallskip
\]
\end{enonce}

\reponse{$ CE^2$}

\begin{corrige}
Il faut intégrer la puissance $\mathcal{P}_E(t)$ fournie par la source sur toute la durée de la charge du condensateur, c'est à dire de $t=0$ à $t=+\infty$. On a donc
\[
\mathcal{E}_E=\int_{t=0}^{t=+\infty} \mathcal{P}_E(t) \, \d{}t = \int_{t=0}^{t=+\infty} \frac{CE^2}{\tau}\exp(-t/\tau) \d{}t = (-\tau) \frac{CE^2}{\tau} \Big[\exp(-t/\tau)\Big]_0^{+\infty}= CE^2.
\]
Remarquons que cette expression est homogène à l'énergie contenue dans un condensateur $\frac{1}{2}Cu_C^2$.
\end{corrige}

% ************************************** %


% ^^^^^^^^^^^^^^^^^^^^^^^^^^^^^^^^^^^^^^ %
%               Question                 %
% ^^^^^^^^^^^^^^^^^^^^^^^^^^^^^^^^^^^^^^ %

\begin{enonce}
l'énergie totale $\mathcal{E}_J$ dissipée par effet Joule que l'on calculera grâce à la formule \minusSmallskip
\[
\mathcal{E}_J = \int_0^\infty  \mathcal{P}_J(t) \d t.\minusBigskip\minusSmallskip
\]
\end{enonce}

\reponse{$ \frac{1}{2} CE^2$}

\begin{corrige}
Il faut intégrer la puissance $\mathcal{P}_J(t)$ sur tout le temps de la charge du condensateur, de $t=0$ à $t=+\infty$ : 
\[
\mathcal{E}_J=\int_{t=0}^{t=+\infty} \mathcal{P}_J(t) \, \d{}t  = \int_{t=0}^{t=+\infty} \frac{CE^2}{\tau}\exp(-2t/\tau) \d{}t = \frac{CE^2}{\tau} \left(-\frac{\tau}{2} \right) \Big[\exp(-2t/\tau)\Big]_0^{+\infty} =\frac{1}{2} CE^2.
\]
\end{corrige}

% ************************************** %


% ^^^^^^^^^^^^^^^^^^^^^^^^^^^^^^^^^^^^^^ %
%               Question                 %
% ^^^^^^^^^^^^^^^^^^^^^^^^^^^^^^^^^^^^^^ %


\begin{enonce}
l'énergie totale $\mathcal{E}_C$ fournie au condensateur que l'on calculera grâce à la formule \minusSmallskip
\[
\mathcal{E}_C = \int_0^\infty  \mathcal{P}_C(t) \d t.\minusBigskip\minusSmallskip
\]
\end{enonce}

\reponse{$\frac{1}{2}CE^2 $}

\begin{corrige}
Il faut intégrer la puissance $\mathcal{P}_C(t)$ sur tout le temps de la charge du condensateur, c'est à dire de $t=0$ à $t=+\infty$. On a donc
\[
\mathcal{E}_C=\int_{t=0}^{t=+\infty} \mathcal{P}_C(t) \, \d{}t  = \int_{t=0}^{t=+\infty} \frac{CE^2}{\tau}(\exp(-t/\tau)-\exp(-2t/\tau)) \d{}t. 
\]
On  reconnaît les deux intégrales précédentes donc
\[
\mathcal{E}_C=  (-\tau) \frac{CE^2}{\tau} \Big[\exp(-t/\tau)\Big]_0^{+\infty} - \frac{CE^2}{\tau} \left(-\frac{\tau}{2} \right) \Big[\exp(-2t/\tau)\Big]_0^{+\infty}=\frac{1}{2}CE^2.
\]
Alternativement, on aurait pu effectuer le calcul suivant:
\[
\mathcal{E}_C=\int_{t=0}^{t=+\infty} \mathcal{P}_C(t) \, \d{}t=\int_{t=0}^{t=+\infty} u_C i \, \d{}t =\int_{t=0}^{t=+\infty} u_C \cdot C \diff{u_C}{t} \, \d{}t= \int_{t=0}^{t=+\infty} \d{} \left(\frac{1}{2}Cu_C^2 \right) 
\]
pour trouver
\[
\mathcal{E}_C=\frac{1}{2}C\Big(u_C^2(+\infty)-u_C^2(0)\Big)=\frac{1}{2}CE^2,
\]
qui est le même résultat.

Remarquons que par conservation de l'énergie, cette dernière expression peut s'obtenir en faisant la différence entre les deux précédentes, l'énergie reçue par le condensateur étant égale à l'énergie fournie par la source de tension dont on a retranché l'énergie dissipée dans le conducteur ohmique. C'est un bon moyen de contrôler le résultat.
\end{corrige}

% ************************************** %


% =============================================================================== %
\finEntrainement
% =============================================================================== %



\newpage



% ********************************************************************* %
%                           ENTRAÎNEMENT                                % 
% ********************************************************************* %
% ================ Métadonnées sur l'entraînement ===================== %

\titreEntrainementFacultatif	{Aspects énergétiques du circuit RLC}
\hauteurLargeurCadreReponse		{8mm}{4cm}
\dureeResolutionFacultative		{3} % 1, 2, 3 ou 4
\typeEntrainement				{} % Newton ou QCM ou AN ou vide (exercice "normal")
\nombreColonnesQuestions		{1} % vide, 1, 2, 3, etc.
\avecPlusieursQuestions			{Y} % Y ou N
\initialisationEntrainement
% ===================================================================== %

                                                         %

On considère le montage ci-dessous dans lequel le condensateur est initialement déchargé.


\begin{center}
\subimport{_images/}{fiche_ELC04_circuit_RLC_1_CBD_v2.tex}
\end{center}

\`A $t=0$, on ferme l'interrupteur $K$.

À $t=0^+$, on a $u_C(t=0^+)=0$ et 
$i(t=0^+)=0$. 

En régime permanent, on a $u_C(t\rightarrow +\infty)=E$ et $i(t\rightarrow +\infty)=0$.

\bigskip


% =============================================================================== %
\debutEntrainement
% =============================================================================== %


% ^^^^^^^^^^^^^^^^^^^^^^^^^^^^^^^^^^^^^^ %
%               Question                 %
% ^^^^^^^^^^^^^^^^^^^^^^^^^^^^^^^^^^^^^^ %

\begin{enonce}
Exprimer la puissance instantanée $\mathcal{P}_E(t)$ fournie par la source en fonction de $E$ et de $u_C(t)$.

{\itshape On pourra s'aider de la relation $i(t)=C\diff{u_C}{t}$.} \minusBigskip

\end{enonce}

\reponse{$EC\diff{u_C}{t}$}

\begin{corrige}
La puissance instantanée délivrée par le source de tension s'écrit $\mathcal{P}_E(t)=Ei(t)=EC\diff{u_C}{t}$.
\end{corrige}

% ************************************** %


% ^^^^^^^^^^^^^^^^^^^^^^^^^^^^^^^^^^^^^^ %
%               Question                 %
% ^^^^^^^^^^^^^^^^^^^^^^^^^^^^^^^^^^^^^^ %

\begin{enonce}
Exprimer la puissance instantanée $\mathcal{P}_C(t)$ reçue par le condensateur en fonction de $u_C(t)$ et $C$. 

\end{enonce}

\reponse{$\diff{\big(\frac{1}{2}Cu_C^2(t)\big)}{t}$}

\begin{corrige}
La puissance instantanée reçue par le condensateur s'écrit 
\[
\mathcal{P}_C(t)=u_C(t)i(t)=u_C(t)C\diff{u_C}{t}=\diff{\big(\frac{1}{2}Cu_C^2(t)\big)}{t}.
\]
\end{corrige}

% ************************************** %


% ^^^^^^^^^^^^^^^^^^^^^^^^^^^^^^^^^^^^^^ %
%               Question                 %
% ^^^^^^^^^^^^^^^^^^^^^^^^^^^^^^^^^^^^^^ %

\begin{enonce}
Exprimer la puissance instantanée $\mathcal{P}_L(t)$ reçue par la bobine en fonction de $i(t)$ et $L$. 

\end{enonce}

\reponse{$\diff{\big(\frac{1}{2}Li^2(t)\big)}{t}$}

\begin{corrige}
La puissance instantanée reçue par la bobine s'écrit $\mathcal{P}_L(t)=u_L(t)i(t)=L\diff{i}{t}i(t)=\diff{\big(\frac{1}{2}Li^2(t)\big)}{t}$.
\end{corrige}

% ************************************** %

\bigskip

En intégrant les expressions des puissances instantanées aux bornes de chaque dipôle, exprimer en fonction des grandeurs introduites:


% ^^^^^^^^^^^^^^^^^^^^^^^^^^^^^^^^^^^^^^ %
%               Question                 %
% ^^^^^^^^^^^^^^^^^^^^^^^^^^^^^^^^^^^^^^ %

\begin{enonce}
L'énergie totale fournie par la source de tension
\end{enonce}

\reponse{$CE^2$}

\begin{corrige}
On intègre la puissance $\mathcal{P}_E(t)$ sur tout le temps de la charge du condensateur, de $t=0$ à $t=+\infty$ : 
\[
\mathcal{E}_E= \int_{t=0}^{t=+\infty} E C \diff{u_C}{t} \d{} t = CE \int_{t=0}^{t=+\infty} \d{} u_C =CE\big( u_C(t=+\infty) - u_C(t=0)\big)= CE^2.
\]
\end{corrige}

% ************************************** %


% ^^^^^^^^^^^^^^^^^^^^^^^^^^^^^^^^^^^^^^ %
%               Question                 %
% ^^^^^^^^^^^^^^^^^^^^^^^^^^^^^^^^^^^^^^ %

\begin{enonce}
L'énergie totale fournie au condensateur
\end{enonce}

\reponse{$\frac{1}{2}CE^2 $}

\begin{corrige}
On intègre la puissance $\mathcal{P}_C(t)$ sur tout le temps de la charge du condensateur, de $t=0$ à $t=+\infty$ :
\[
\mathcal{E}_C=\int_{t=0}^{t=+\infty} \mathcal{P}_C(t) \, \d{} t = \int_{t=0}^{t=+\infty} \d{}\left(\frac{1}{2}Cu_C^2 \right) 
= \frac{1}{2}C\big({u_C}^2(+\infty)-{u_C}^2(0)\big)=\frac{1}{2}CE^2.
\]
\end{corrige}

% ************************************** %


% ^^^^^^^^^^^^^^^^^^^^^^^^^^^^^^^^^^^^^^ %
%               Question                 %
% ^^^^^^^^^^^^^^^^^^^^^^^^^^^^^^^^^^^^^^ %

\begin{enonce}
L'énergie totale fournie à la bobine
\end{enonce}

\reponse{$0$}

\begin{corrige}
On intègre la puissance $\mathcal{P}_L(t)$ sur tout le temps de la charge du condensateur, de $t=0$ à $t=+\infty$ :
\[
\mathcal{E}_L=\int_{t=0}^{t=+\infty} \mathcal{P}_L(t) \, \d t=\int_{t=0}^{t=+\infty} \ d\left(\frac{1}{2}Li^2 \right) %\, dt 
=\frac{1}{2}L(i^2(+\infty)-i^2(0))=0.
\]
\end{corrige}

% ************************************** %





% ^^^^^^^^^^^^^^^^^^^^^^^^^^^^^^^^^^^^^^ %
%               Question                 %
% ^^^^^^^^^^^^^^^^^^^^^^^^^^^^^^^^^^^^^^ %

\begin{enonce}
En exploitant les résultats précédents, exprimer l'énergie totale dissipée par effet Joule.

\end{enonce}

\reponse{$ \frac{1}{2} CE^2$}

\begin{corrige}
Il faudrait intégrer la puissance dissipée par effet Joule $\mathcal{P}_J(t)=Ri^2(t)$ sur tout le temps de la charge du condensateur, de $t=0$ à $t=+\infty$.  Mais, on n'a pas accès à l'expression de $i(t)$. On peut alors malgré tout se servir 
de la conservation de l'énergie:
\[
\mathcal{E}_J=\mathcal{E}_E-\mathcal{E}_C-\mathcal{E}_L= CE^2-\frac{1}{2} CE^2-0=\frac{1}{2} CE^2.
\]
\end{corrige}

% ************************************** %


% =============================================================================== %
\finEntrainement
% =============================================================================== %




\newpage



% •••••••••••••••••••••••••••••••••••••••••••••••••••••••••••••••••••••••••••• %
% •••••••••••••••••••••••••••••••••••••••••••••••••••••••••••••••••••••••••••• %
\sectionFicheEntrainement{Bilan d'énergie en régime sinusoïdal forcé}
% •••••••••••••••••••••••••••••••••••••••••••••••••••••••••••••••••••••••••••• %
% •••••••••••••••••••••••••••••••••••••••••••••••••••••••••••••••••••••••••••• %



%
%% ********************************************************************* %
%%                           ENTRAÎNEMENT                                % 
%% ********************************************************************* %
%% ================ Métadonnées sur l'entraînement ===================== %
%
%\titreEntrainementFacultatif	{Circuit RLC en régime sinusoïdal forcé}
%\hauteurLargeurCadreReponse		{8mm}{4cm}
%\dureeResolutionFacultative		{3} % 1, 2, 3 ou 4
%\typeEntrainement				{} % Newton ou QCM ou AN ou vide (exercice "normal")
%\nombreColonnesQuestions		{1} % vide, 1, 2, 3, etc.
%\avecPlusieursQuestions			{Y} % Y ou N
%\initialisationEntrainement
%% ===================================================================== %
%
%Un circuit RLC série de pulsation propre $\omega_0$ et de résistance $R$ est excité à sa pulsation de résonance par un générateur de tension $e(t)=E\cos(\omega_0 t)$.
%Dans ces conditions, l'expression de l'intensité $i(t)$ est  $i(t)=I_0\cos (\omega_0 t +\varphi)$. On rappelle que la pulsation de résonance vérifie $\omega_0^2=\frac{1}{LC}$ et le facteur de qualité du circuit $Q=\frac{1}{RC\omega_0}$.
%
%Exprimer, en fonction des grandeurs introduites:
%
%% =============================================================================== %
%\debutEntrainement
%% =============================================================================== %
%
%
%
%
%% ^^^^^^^^^^^^^^^^^^^^^^^^^^^^^^^^^^^^^^ %
%%               Question                 %
%% ^^^^^^^^^^^^^^^^^^^^^^^^^^^^^^^^^^^^^^ %
%
%\begin{enonce}
%l'expression de la tension aux bornes du condensateur $u_C(t)$. On utilisera pour cela la relation $i(t)=C\frac{du_C}{dt}$ aux bornes du condensateur.
%\end{enonce}
%
%\reponse{$\frac{I_0}{C \omega_0} \sin (\omega_0 t +\varphi)$}
%
%\begin{corrige}
%Par définition, dans un circuit RLC série, $i(t)=C\diff{u_C}{t}$, donc en régime sinusoïdal forcé
%$$
%i(t)=C\diff{u_C}{t} \Rightarrow u_C(t)=\frac{1}{C}\int i(t) dt = \frac{I_0}{C \omega_0} \sin (\omega_0 t +\varphi)
%$$
%\end{corrige}
%
%% ************************************** %
%
%
%% ^^^^^^^^^^^^^^^^^^^^^^^^^^^^^^^^^^^^^^ %
%%               Question                 %
%% ^^^^^^^^^^^^^^^^^^^^^^^^^^^^^^^^^^^^^^ %
%
%\begin{enonce}
%l'énergie $\mathcal{E}_C(t)= \frac{1}{2}Cu_C^2(t)$ stockée dans le condensateur.
%\end{enonce}
%
%\reponse{$\frac{1}{2}\frac{I_0^2}{C\omega_0^2} \sin^2 (\omega_0 t +\varphi)$}
%
%\begin{corrige}
% L'énergie stockée dans le condensateur vaut
% $$\mathcal{E}_C(t)=\frac{1}{2}Cu_C^2=\frac{1}{2}C\frac{I_0^2}{C^2 \omega_0^2} \sin^2 (\omega_0 t +\varphi)=\frac{1}{2}\frac{I_0^2}{C\omega_0^2} \sin^2 (\omega_0 t +\varphi)$$
%\end{corrige}
%
%% ************************************** %
%
%
%% ^^^^^^^^^^^^^^^^^^^^^^^^^^^^^^^^^^^^^^ %
%%               Question                 %
%% ^^^^^^^^^^^^^^^^^^^^^^^^^^^^^^^^^^^^^^ %
%
%\begin{enonce}
%l'énergie $\mathcal{E}_L(t)=\frac{1}{2}Li^2(t)$ stockée dans la bobine.
%\end{enonce}
%
%\reponse{$\frac{1}{2}LI_0^2\cos^2 (\omega_0 t +\varphi)$}
%
%\begin{corrige}
% L'énergie stockée dans la bobine vaut
% $$\mathcal{E}_L(t)=\frac{1}{2}Li^2=\frac{1}{2}LI_0^2\cos^2 (\omega_0 t +\varphi)$$
%\end{corrige}
%
%% ************************************** %
%
%
%% ^^^^^^^^^^^^^^^^^^^^^^^^^^^^^^^^^^^^^^ %
%%               Question                 %
%% ^^^^^^^^^^^^^^^^^^^^^^^^^^^^^^^^^^^^^^ %
%
%\begin{enonce}
%l'énergie totale $\mathcal{E}_{tot}=\mathcal{E}_L(t)+\mathcal{E}_C(t)$ stockée dans l'oscillateur en fonction de $C$, $I_0$ et $\omega_0$. 
%
%\end{enonce}
%
%\reponse{$\frac{1}{2C\omega_0^2} I_0^2$}
%
%\begin{corrige}
%L'énergie totale de l'oscillateur est constituée de l'énergie de la bobine et de celle du condensateur
%$$
%\mathcal{E}_{tot}=\frac{1}{2}\left( L I_0^2\cos^2 (\omega_0 t +\varphi)+ \frac{I_0^2}{C \omega_0^2} \sin^2 (\omega_0 t +\varphi)\right)
%$$
%que l'on peut simplifier en utilisant l'indication de l'énoncé
%$$
%\mathcal{E}_{tot}=\frac{1}{2}\left( \frac{1}{C\omega_0^2} I_0^2\cos^2 (\omega_0 t +\varphi)+ \frac{I_0^2}{C \omega_0^2} \sin^2 (\omega_0 t +\varphi)\right)=\frac{1}{2}\frac{1}{C\omega_0^2} I_0^2(\cos^2 (\omega_0 t +\varphi)+\sin^2 (\omega_0 t +\varphi))
%$$
%et finalement
%$$
%\mathcal{E}_{tot}=\frac{1}{2C\omega_0^2} I_0^2
%$$
%\end{corrige}
%
%% ************************************** %
%
%
%% ^^^^^^^^^^^^^^^^^^^^^^^^^^^^^^^^^^^^^^ %
%%               Question                 %
%% ^^^^^^^^^^^^^^^^^^^^^^^^^^^^^^^^^^^^^^ %
%
%\begin{enonce}
%la puissance $\mathcal{P}_J(t)$ dissipée dans le conducteur ohmique.
% \end{enonce}
%
%\reponse{$RI_0^2\cos^2 (\omega_0 t +\varphi)$}
%
%\begin{corrige}
%Par définition, $\mathcal{P}_J(t)=Ri^2$ donc
%$$
%\mathcal{P}_J(t)=RI_0^2\cos^2 (\omega_0 t +\varphi)
%$$
%\end{corrige}
%
%% ************************************** %
%
%
%% ^^^^^^^^^^^^^^^^^^^^^^^^^^^^^^^^^^^^^^ %
%%               Question                 %
%% ^^^^^^^^^^^^^^^^^^^^^^^^^^^^^^^^^^^^^^ %
%
%\begin{enonce}
%l'énergie $\Delta \mathcal{E}=\int_0^{T_0} \mathcal{P}_J(t)dt$ dissipée dans le conducteur ohmique pendant une période \\ $T_0=\frac{2\pi}{\omega_0}$ de l'oscillateur en fonction de $R$, $I_0$ et $\omega _0$.
%\end{enonce}
%
%\reponse{$RI_0^2\frac{\pi}{\omega_0}$}
%
%\begin{corrige}
%On calcule
%$$
%\Delta \mathcal{E}=\int_0^{T_0} RI_0^2\cos^2 (\omega_0 t +\varphi) \d{}t 
%$$
%On linéarise le cosinus avec la formule $\cos^2 \alpha= \frac{1}{2} \left(1+\cos 2\alpha\right)$, ce qui donne
%$$
%\Delta \mathcal{E}=RI_0^2  \int_0^{T_0}  \frac{1}{2} \left[1+\cos\left(2\omega_0 t +2\varphi\right)\right] \d{}t = \frac{RI_0^2}{2} \left[\int_0^{T_0} \d{}t  + \int_0^{T_0} \cos\left(2\omega_0 t +2\varphi\right) \d{}t  \right]
%$$
%La première intégrale donne $T_0$, la deuxième partie s'intègre en $-\sin$, il reste
%$$
%\Delta \mathcal{E}=\frac{RI_0^2}{2} \left(T_0-\left(\frac{1}{2\omega_0}\right)\left[\sin(2\omega_0 t +2\varphi)\right]_0^{T_0}  \right)
%$$
%Par définition de la périodicité du sinus, le terme $\left[\sin(2\omega_0 t +2\varphi)\right]_0^{T_0}$ est nul et donc:
%$$
%\Delta \mathcal{E}=\frac{RI_0^2}{2} T_0=\frac{RI_0^2}{2} \frac{2\pi}{\omega_0}
%$$
%\end{corrige}
%
%% ************************************** %
%
%
%% ^^^^^^^^^^^^^^^^^^^^^^^^^^^^^^^^^^^^^^ %
%%               Question                 %
%% ^^^^^^^^^^^^^^^^^^^^^^^^^^^^^^^^^^^^^^ %
%
%\begin{enonce}
%le rapport $\frac{\mathcal{E}_{tot}}{\Delta \mathcal{E}}$ en fonction du facteur de qualité $Q$.
%\end{enonce}
%
%\reponse{$\frac{Q}{2\pi}$}
%
%\begin{corrige}
%On calcule la quantité demandée
%$$
%\frac{\mathcal{E}_{tot}}{\Delta \mathcal{E}}=\frac{1}{2C\omega_0^2} I_0^2 \cdot \frac{2} {RI_0^2} \frac{\omega_0}{2\pi}= \frac{1}{2\pi RC \omega_0}
%$$
%soit finalement
%$$
%\frac{\mathcal{E}_{tot}}{\Delta \mathcal{E}}=\frac{Q}{2\pi}
%$$
%\end{corrige}
%
%% ************************************** %
%
%
%% =============================================================================== %
%\finEntrainement
%% =============================================================================== %
%
%




% ********************************************************************* %
%                           ENTRAÎNEMENT                                % 
% ********************************************************************* %
% ================ Métadonnées sur l'entraînement ===================== %

\titreEntrainementFacultatif	{Adaptation d'impédance}
\hauteurLargeurCadreReponse		{8mm}{3cm}
\dureeResolutionFacultative		{3} % 1, 2, 3 ou 4
\typeEntrainement				{Newton} % Newton ou QCM ou AN ou vide (exercice "normal")
\nombreColonnesQuestions		{1} % vide, 1, 2, 3, etc.
\avecPlusieursQuestions			{Y} % Y ou N
\initialisationEntrainement
% ===================================================================== %

% mmmmmmmmmmmmmmmmmmmmmmmmmmmmmmmmmmmmmmmmmmmmmmmmmmmmmmmmm %
% 		  Insertion d'une image en regard d'un texte
% mmmmmmmmmmmmmmmmmmmmmmmmmmmmmmmmmmmmmmmmmmmmmmmmmmmmmmmmm %
% --------------------------------------------------------- %
\pourcentageDeLaPartieAGauche   {0.55}
% --------------------------------------------------------- %
%                                                           %
% -------------------- Partie à gauche -------------------- %
%                                                           %
                                \initialisationPartieGauche % 
%                                                           %
%                                                           %

On considère un dipôle d'impédance $\underline{Z_u}$ branché aux bornes d'un générateur de fém $\underline{e_G}(t)$ et d'impédance interne $\underline{Z_G}$. 

On notera: $\underline{Z_u} = R_u +\jC X_u$ et $\underline{Z_G} = R_G +\jC X_G$.

Le dipôle $\underline{Z_u}$ est traversé par le courant d'intensité $i(t)$.

On écri en notation complexe, 
\[
\underline{e_G}  = E\sqrt{2} \eEuler^{\jC \omega t}
\quad\text{et}\quad
\underline{i}  = I\sqrt{2} \eEuler^{\jC (\omega t + \varphi)}.
\]

% --------------------------------------------------------- %
%                                                           %
% -------------------- Partie à droite -------------------- %
%                                                           %
                                \initialisationPartieDroite %
%                                                           %
%                                                           %
\begin{center}
\subimport{_images/}{fiche_adapt_RSF_Complex.tex}
\end{center}
% --------------------------------------------------------- %
                               \finalisationDuPartageDePage %					
% --------------------------------------------------------- %




% =============================================================================== %
\debutEntrainement
% =============================================================================== %


La puissance moyenne reçue par l'impédance $\underline{Z_{u}}$ vaut
\[
\mathcal{P}_{m}=\dfrac{1}{2}\operatorname{Re}\left(\underline{Z_{u}}\times \underline{i}\times\underline{i}^{\star}\right).
\] 




% ^^^^^^^^^^^^^^^^^^^^^^^^^^^^^^^^^^^^^^ %
%               Question                 %
% ^^^^^^^^^^^^^^^^^^^^^^^^^^^^^^^^^^^^^^ %

\begin{enonce}
Exprimer la puissance $\mathcal{P}_{m}$ en fonction de $I$ et de $R_u$
\end{enonce}

\reponse{$R_{u}I^{2}$}

\begin{corrige}
On a 
\[
\mathcal{P}_{m}  =\dfrac{1}{2}\operatorname{Re}\left(\left(R_{u}+jX_{u}\right)I\sqrt{2}\eEuler^{\jC \left(\omega t+\varphi\right)}\cdot I\sqrt{2}\eEuler^{-j\left(\omega t+\varphi\right)}\right)=\operatorname{Re}\left(R_{u}+jX_{u}\right)I^{2}
  =R_{u}I^{2}.
  \]
\end{corrige}

% ************************************** %


% ^^^^^^^^^^^^^^^^^^^^^^^^^^^^^^^^^^^^^^ %
%               Question                 %
% ^^^^^^^^^^^^^^^^^^^^^^^^^^^^^^^^^^^^^^ %

\begin{enonce}
Grâce à une loi des mailles, exprimer $I$ en fonction de $E$, de $R_{G}$, $R_{u}$ et $X_{G}$, $X_{u}$.

\end{enonce}

\reponse{$\dfrac{E}{\sqrt{\left(R_{G}+R_{u}\right)^{2}+\left(X_{G}+X_{u}\right)^{2}}}$}

\begin{corrige}
La loi des mailles donne:
\begin{align*}
\underline{e_{G}} & =\left(\underline{Z_{G}}+\underline{Z_{u}}\right)\underline{i}\\
\text{donc } E\sqrt{2}\eEuler^{\jC \omega t} & =\left[R_{G}+R_{u}+j\left(X_{G}+X_{u}\right)\right]I\sqrt{2}\eEuler^{\jC \left(\omega t+\varphi\right)}\\
\text{donc } E & =\left[R_{G}+R_{u}+j\left(X_{G}+X_{u}\right)\right]I\eEuler^{\jC \varphi}.
\end{align*}
En prenant le module, on obtient
\[
I=\dfrac{E}{\sqrt{\left(R_{G}+R_{u}\right)^{2}+\left(X_{G}+X_{u}\right)^{2}}}.
\]
\end{corrige}

% ************************************** %


% =============================================================================== %
\pauseEntrainement
% =============================================================================== %


Des résultats précédents, on en déduit l'expression de $\mathcal{P}_{m}$ en fonction de $E$:
\[
\mathcal{P}_{m}=\dfrac{R_{u}E^2}{\left(R_{G}+R_{u}\right)^{2}+\left(X_{G}+X_{u}\right)^{2}}
\]


On cherche à déterminer les conditions sur $R_u$ et $X_u$ pour que $\mathcal{P}_m$ soit maximale. On dit alors qu'il y a \emph{adaptation d'impédance}.



% =============================================================================== %
\repriseEntrainement
% =============================================================================== %


% ^^^^^^^^^^^^^^^^^^^^^^^^^^^^^^^^^^^^^^ %
%               Question                 %
% ^^^^^^^^^^^^^^^^^^^^^^^^^^^^^^^^^^^^^^ %

\begin{enonce}
Calculer la dérivée de $\mathcal{P}_m$ par rapport à $X_u$
\end{enonce}
\reponse{\small$-R_uE^2\frac{2(X_G+X_u)}{\Big((R_G+R_u)^2+(X_G+X_u)^2\Big)^2}$}

\begin{corrigeNewpage}
En reportant l'expression de $I$ obtenue dans celle de $\mathcal{P}_m$ on retrouve l'expression donnée dans l'énoncé : 
\[
\mathcal{P}_m=\frac{R_uE^2}{(R_G+ R_u)^2+(X_G+X_u)^2}.
\]
La fonction dont il faut calculer la dérivée est du type \(\mathcal{P}_m(X_u)=\frac{1}{f(X_u)}\). La dérivée sera donc du type 
\[
\frac{\partial \mathcal{P}_m}{\partial X_u} = -\frac{f'(X_u)}{(f(X_u))^2}.
\]
 Finalement, on calcule : 
\[
\frac{\partial \mathcal{P}_m}{\partial X_u}=-R_uE^2\frac{2(X_G+X_u)}{\Big((R_G+R_u)^2+(X_G+X_u)^2\Big)^2}.
\]
\end{corrigeNewpage}

% ************************************** %


% ^^^^^^^^^^^^^^^^^^^^^^^^^^^^^^^^^^^^^^ %
%               Question                 %
% ^^^^^^^^^^^^^^^^^^^^^^^^^^^^^^^^^^^^^^ %

\begin{enonce}
Calculer la dérivée de $\mathcal{P}_m$ par rapport à $R_u$
\end{enonce}

\reponse{$E^2\frac{(R_G^2-R_u^2)+(X_G+X_u)^2}{\Big((R_G+R_u)^2+(X_G+X_u)^2\Big)^2}$}

\begin{corrige}
La fonction dont il faut calculer la dérivée est du type \(\mathcal{P}_m(R_u)=\frac{f(R_u)}{g(R_u)}\), la dérivée sera donc du type 
\[
\frac{\partial \mathcal{P}_m}{\partial R_u} = \frac{f'(R_u)g(R_u)-f(R_u)g'(R_u)}{(g(R_u))^2}.
\]
 Ainsi, on calcule
\begin{align*}
\frac{\partial \mathcal{P}_m}{\partial R_u}& =  E^2\frac{(R_G+R_u)^2+(X_G+X_u)^2-2R_u(R_G+R_u)}{\Big((R_G+R_u)^2+(X_G+X_u)^2\Big)^2}\\
 & =  E^2\frac{R_G^2+R_u^2+2R_GR_u+(X_G+X_u)^2-2R_u^2-2R_uR_G}{\Big((R_G+R_u)^2+(X_G+X_u)^2\Big)^2}\\
   & =  E^2\frac{(R_G^2-R_u^2)+(X_G+X_u)^2}{\Big((R_G+R_u)^2+(X_G+X_u)^2\Big)^2}.
\end{align*}
\end{corrige}

% ************************************** %


% ^^^^^^^^^^^^^^^^^^^^^^^^^^^^^^^^^^^^^^ %
%               Question                 %
% ^^^^^^^^^^^^^^^^^^^^^^^^^^^^^^^^^^^^^^ %

\begin{enonce}
Choisir parmi les quatre propositions suivantes quelle est la condition pour que $\mathcal{P}_m$ soit maximale:
	\begin{listeQCM2Colonnes}
	\item $X_u=-X_G$ et $R_u=-R_G$
	\item $X_u=X_G$ et $R_u=-R_G$
	\item $X_u=-X_G$ et $R_u=R_G$
	\item $X_u=X_G$ et $R_u=R_G$
	\end{listeQCM2Colonnes}

\end{enonce}

\reponse{\reponseC{}}

\begin{corrige}
On cherche pour quelles valeurs de $R_u$ et $X_u$ les deux dérivées partielles de $\mathcal{P}_m$ sont nulles.

On a $\frac{\partial \mathcal{P}_m}{\partial X_u}=0$ pour $X_u+X_G=0$ soit $X_u=-X_G$.

On aura alors $\frac{\partial \mathcal{P}_m}{\partial R_u}=E^2\frac{(R_G^2-R_u^2)}{\Big((R_G+R_u)^2+(X_G+X_u)^2\Big)^2}$. Alors, on a  $\frac{\partial \mathcal{P}_m}{\partial R_u}=0$ pour $R_G=R_u$.

{\itshape Mathématiquement, on pourrait avoir comme solution $R_G=R_u$ ou $R_G=-R_u$. Ainsi, la solution \reponseA{} pourrait  aussi être considérée comme correcte. Mais, en physique, on a nécessairement $R_G\geq 0$ et $R_u\geq 0$.}
\end{corrige}

% ************************************** %

% =============================================================================== %
\finEntrainement
% =============================================================================== %








% ---------------------------------------------------------------------------- %
%                       Affichage des réponses mélangées                       %
% ---------------------------------------------------------------------------- %
\afficheReponsesMelangees
% ---------------------------------------------------------------------------- %

%%%%%%%%%%%%%%%%%%%%%%%%%%%%%%%%%%%%%%%%%%%%%%%%%%%%%%%%%%%%%%%%%%%%%%%%%%%%%%%%%%%%%%%%%%%%%%
\finFicheEntrainement                                                                        %
%%%%%%%%%%%%%%%%%%%%%%%%%%%%%%%%%%%%%%%%%%%%%%%%%%%%%%%%%%%%%%%%%%%%%%%%%%%%%%%%%%%%%%%%%%%%%%


% ======================================================= % 
% ============== Gestion de la compilation ============== %
% ======================================================= % 
\ifdefined\mainIsLoaded\else
\printReponsesEtCorriges
\end{document}\fi
% ======================================================= % 
% ======================================================= % 